\documentclass[11pt,a4paper]{article}
\usepackage[T1]{fontenc}
\usepackage{lmodern,amsmath,amssymb,amsthm}
\usepackage[margin=26mm]{geometry}
\usepackage[hypertexnames=false,hidelinks,bookmarksnumbered=true]{hyperref}
\hypersetup{pdftitle={Hausdorff dimension of Weierstrass cosine graphs for real bases},
  pdfauthor={Timofei Iuzhakov}, pdfsubject={Version 1.0.1. Proposed proof of the Hausdorff dimension formula for the classical Weierstrass function at all real bases and with prescribed phases},
  pdfkeywords={classical Weierstrass function, Hausdorff dimension, graph of the Weierstrass function, non-integer bases, prescribed phases, exact-dimensional measures}}
\setlength{\emergencystretch}{2em}
\newtheorem{theorem}{Theorem}[section]
\newtheorem{proposition}[theorem]{Proposition}
\newtheorem{lemma}[theorem]{Lemma}
\newcommand{\T}{\mathbb T}
\newcommand{\Leb}{\operatorname{Leb}}
\newcommand{\E}{\mathbb E}
\title{Hausdorff dimension of Weierstrass cosine graphs for real bases}
\author{Timofei Iuzhakov}
\date{\normalsize Version 1.0.1 --- manuscript for external review --- 19 November 2026\\[2pt]
  {\small \href{https://github.com/ppadfoot/weierstrass-real-bases/releases/tag/v1.0.1}{\underline{GitHub release and source}}\quad\textperiodcentered\quad\href{https://doi.org/10.5281/zenodo.24378509}{\underline{Zenodo DOI 10.5281/zenodo.24378509}}}}
\begin{document}
\maketitle
\begin{abstract}
For real $b>1$ and $b^{-1}<\lambda<1$, the conjectured Hausdorff dimension
of the graph of the classical Weierstrass function
$W(x)=\sum_{j\ge0}\lambda^j\cos(2\pi b^j x)$ is
$2+\log\lambda/\log b$. We present a proposed proof for every real
base. The argument also gives this dimension for the graph with any
prescribed sequence of phase offsets and for its natural measure on
each bounded interval of positive length. The lower bound is reduced
to an estimate for small increments of a two-sided series. At
transcendental bases the associated phases are independent under Haar
measure. At algebraic bases the proof uses either the finite number of
predecessors in the phase group or a contracting conjugate. In the
latter case reciprocal changes of base and transverse differentiation
propagate the increment estimate between amplitudes.
\end{abstract}

\section{The graph dimension problem}\label{sec:intro}
Let $b>1$, $b^{-1}<\lambda<1$, $\T=\mathbb R/\mathbb Z$ and
$\alpha=-\log\lambda/\log b\in(0,1)$. For prescribed phases
$\varphi=(\varphi_j)_{j\ge0}\in\T^{\{0,1,\ldots\}}$, write
\[
 W_\varphi(x)=\sum_{j=0}^{\infty}\lambda^j
              \cos\bigl(2\pi(b^jx+\varphi_j)\bigr),\qquad x\in\mathbb R.
\]
Write $W=W_\varphi$ when the phases are fixed. The series converges
uniformly. Splitting the sum at the frequency nearest $1/|u|$ bounds
$|W(x+u)-W(x)|$ by a constant times $|u|^\alpha$. The resulting
H\"older cover gives $2-\alpha$ as an upper bound for graph dimension.
The question is whether this upper bound is attained.

\begin{theorem}[Graph and graph measure]\label{thm:main}
For every real $b>1$, every $b^{-1}<\lambda<1$, every prescribed
phase sequence $\varphi$, and every interval $I\subset\mathbb R$ of
positive length,
\[
 \dim_{\mathrm H}\{(x,W_\varphi(x)):x\in I\}=2-\alpha
                    =2+\frac{\log\lambda}{\log b}.
\]
If $I$ is bounded, let $\mu_I$ be the image of normalized Lebesgue
measure on $I$ under $x\mapsto(x,W_\varphi(x))$. This measure is
exact dimensional. For $\mu_I$-almost every $z\in\mathbb R^2$,
\[
 \lim_{r\downarrow0}\frac{\log\mu_I(B(z,r))}{\log r}=2-\alpha,
\]
where $B(z,r)$ is the Euclidean ball of radius $r$.
\end{theorem}

\subsection{The conjecture and earlier results}\label{sec:history}
The classical Weierstrass function is the case $\varphi_j=0$ for all $j$.
Mauldin and Williams record the real-base cosine conjecture on p.~794
of \cite{mw}. Their Theorem~1 identifies the graph dimensions of the
one-sided and two-sided series by addition of a Lipschitz function.
For fixed $\alpha$ they obtain a lower bound $2-\alpha-C/\log b$ for
sufficiently large $b$. Hunt obtains $2-\alpha$ for each real base
for almost every sequence of independent uniform phases \cite{hunt}.
In Theorem~\ref{thm:main} the phases may be prescribed in advance,
including the sequence consisting entirely of zeros.

For integer bases, Bara\'nski, B\'ar\'any and Romanowska prove
the expected dimension throughout a range $\lambda_b<\lambda<1$
\cite{bbr}. Keller gives a proof of this range by estimating the
distribution of slopes at successive scales \cite{keller}. Shen
establishes the classical cosine formula for every integer $b\ge2$
throughout $1/b<\lambda<1$ \cite{shen}. In this integer-base model,
Ledrappier's criterion gives the natural graph measure dimension
$2-\alpha$ when the slope distribution has dimension one for
Lebesgue-almost every base point \cite{ledrappier}, as stated in
\cite[Theorem~2.2 and Corollary~2.3]{bbr}.
For a phase shared by all terms, Ruan and Zhang prove the graph
formula for every integer $b\ge7$, every $1/b<\lambda<1$ and every
value of that phase \cite[Theorem~1.1]{rz}. Ren and Shen prove a graph
dimension dichotomy for a fixed real-analytic periodic function at an
integer base \cite{rs}. Theorem~F of Li, Pan, Tong and Xu gives the
corresponding dichotomy for a fixed $C^r$ periodic function, $r>1$,
also for integer bases \cite{lptx}. For conditional slope measures,
Ren's dimension theorem for analytic skew products overlaps the
integer-base case of Theorem~\ref{thm:algebraic-slope-dimension}
\cite{ren}.

\subsection{Proof by arithmetic type}\label{sec:routes}
The common estimate concerns the amount of time a rescaled increment
spends within distance $r$ of zero. The assertion $\mathcal S_s$,
defined in \eqref{eq:anchor}, controls that time at each exponent
strictly below $s$ under Haar measure on a coset of the phase group.
Uniformity over cosets is what permits an arbitrary deterministic
phase sequence. Proposition~\ref{prop:graph} converts $\mathcal S_s$
into graph dimension at least $1+s(1-\alpha)$. Thus it suffices
to establish $\mathcal S_1$ for the unweighted cosine series.

\paragraph{Transcendental bases.}
Let $H_b$ be the closure of $x\mapsto (b^j x\bmod 1)_{j\in\mathbb Z}$.
For any nonzero finite integer vector $(h_j)$, the sum $\sum_j h_j b^j$
is nonzero when $b$ is transcendental: multiplying the Laurent
polynomial by a suitable power of its variable gives a nonzero integer
polynomial. No nontrivial character annihilates $H_b$, so
$H_b=\T^{\mathbb Z}$. Its Haar coordinates $z_j$ are therefore independent
and uniform. We apply Hunt's three-term convolution estimate
\cite[pp.~797--798]{hunt} to these Haar increments. For each $t\ne0$ choose $k\in\mathbb Z$ with
$1/(4b)<b^k|t|\le1/4$. In the unweighted two-sided increment
\eqref{eq:increment}, the independent terms with indices $k,k-1,k-2$
have arcsine laws. The amplitude of each is
$2\lambda^j|\sin(\pi b^jt)|\asymp_b |t|^\alpha$.
Rescaling gives each density an $L^{3/2}$ norm
$O(|t|^{-\alpha/3})$. Young's convolution inequality gives their
sum a density bounded by $C|t|^{-\alpha}$, a bound preserved under
convolution with the law of the remaining independent terms.
Integration on $[-T,T]$ bounds the expected time within $Ar$ of
zero by $C_{A,T}r$, since $\alpha<1$. Markov's inequality gives
$\mathcal S_1$ for the unweighted cosine. Equivalently for the graph
argument, the cutoff $B_q$ of Appendix~\ref{app:graph} satisfies
$\int B_q\,dm_{H_b}\le C\lambda^q$. Jensen's inequality yields
$\int -\log\min(1,B_q)\,dm_{H_b}\ge q\log(1/\lambda)-\log C$.
For each prescribed sequence of phases the character estimate in that
appendix identifies the limit on its scale windows with this same
product Haar law for Lebesgue-almost every physical starting point.

\paragraph{Algebraic bases with sufficiently many predecessors.}
Let $P$ be a primitive irreducible integer polynomial of $b$, let
$d=\deg P$, and write $p=|P(0)|$, $A=|\operatorname{lc}(P)|$.
Suppose $p\ge b\lambda$. Then $p\ge2$. Given a central block of
$d$ phases, the preceding phase has $p$ equiprobable choices.
Theorem~\ref{thm:algebraic-slope-dimension} gives the dimension of
the resulting conditional slope measure. For
$\gamma=(b\lambda)^{-1}$ its dimension is
$\min(1,\log p/\log(1/\gamma))=1$, including $p=b\lambda$.
Estimate~\eqref{eq:uniform-slope-information} and
Section~\ref{sec:secant} give $\mathcal S_1$ uniformly over phase
cosets.

\paragraph{Algebraic bases with a contracting conjugate.}
If $p<b\lambda$, the other conjugates $b_2,\ldots,b_d$ satisfy
\[
 \prod_{i=2}^{d}|b_i|=\frac{p}{Ab}<\frac{\lambda}{A}<1.
\]
The empty product for degree one is one, so this case has $d\ge2$.
Since $P(0)\ne0$, at least one conjugate $C$ satisfies $0<|C|<1$.
Use $B=b$ and compare
displacements at $B$ with displacements along a real line at $C$.
Proposition~\ref{prop:two} gives $\mathcal S_1$ for these two
embeddings. Its proof distinguishes the expansion rate
$L=\log|B|$ and the contraction rate $M=-\log|C|$.
For $L=M$, Section~\ref{sec:equal} combines a strip estimate with
the initial exponent $1/2$ of Section~\ref{sec:ray-half-seed}.
For $L\ne M$, Proposition~\ref{prop:transfer} changes the order
of the embeddings, and \eqref{eq:descent} increases the exponent
by transverse differentiation. The recursion in
Section~\ref{sec:closure} combines these steps with the critical
estimate of Section~\ref{sec:weighted-critical-ray}.
For example, at the golden ratio $b=(1+\sqrt5)/2$ and
$\lambda=b^{-1/2}$, one has $p=1<b\lambda$ and $C=-1/b$.
The equal-rate argument gives the value $3/2$.
These cases exhaust the real bases and give Theorem~\ref{thm:main}
by Proposition~\ref{prop:graph}.

\section{Phase groups and reciprocal transfer}\label{sec:embeddings}
Let $B,C$ be conjugate algebraic numbers, $b=|B|=e^L>1$,
$\rho=|C|=e^{-M}<1$. Here and on reversing the pair, at most one of $B,C$ is
nonreal. Put $E_D=\mathbb R$ for a real $D$ and $E_D=\mathbb C$
otherwise. Write
\[
 v_D(e)_j=\operatorname{Re}(eD^j)\pmod1,\qquad
 p_j(v,D)=\operatorname{Re}\big(v(D/|D|)^j\big).
\]
Use the compact subgroup $H_B=\overline{v_B(E_B)}$ of $\T^{\mathbb Z}$, with closure in the product topology. It contains
$v_C(E_C)$, as the annihilator consists of the Laurent polynomials
vanishing at $B$. Reversal $Rz=(z_{-j})_j$ maps $H_B$ onto
$H_{1/C}$, by conjugation and character separation. Write $(Sz)_j=z_{j+1}$ for the shift. The multiplier of the reversed
group is the field element $1/C$ with its sign or complex argument.

Test real displacements $t$ along fixed unit rays $v\in E_B$ and
$w\in E_C$. At a real embedding we permit either sign. For the unique
nonreal embedding, if its angle $\vartheta$ is an irrational multiple
of $\pi$, permit arbitrary unit rays. Otherwise fix a unit $v_0$
with $\operatorname{Re}(v_0e^{ij\vartheta})\ne0$ for all $j$;
permit the finite orbits under rotation by $\vartheta$, conjugation
and sign of $v_0$. These choices work simultaneously for reciprocals.
Write $c=(c_j)$ for any finite product of factors $p_j(v,D)$ from
this family and their reversals, and optionally a constant sign.
Such families with at most $k$ factors are compact, and invariant under
translations of the subscript and reversal. For each fixed $k$ there
are $c_k>0,N_k<\infty$ such that every $N_k$ consecutive indices
contain one where $|c_j|\ge c_k$, uniformly in the family. For a
finite rotation group each permitted factor is bounded away from
zero. For irrational rotation write the complex factors as
$\cos(x+\phi)$, taking absolute values to discard real signs.
Each resulting trigonometric polynomial is nonzero, since its
highest Fourier coefficient has modulus $2^{-k'}$ for $k'\le k$
factors. Compactness and the density of the rotation orbit establish
the uniform finite-window assertion.

Write $\lambda=e^{-a}\in(1/b,1)$ and define the scalar rough function
\begin{equation}\label{eq:increment}
 F_z(t)=\sum_{j\in\mathbb Z}\lambda^jc_j
   \{f(z_j+b^jp_j(v,B)t)-f(z_j)\},\qquad f(y)=\cos(2\pi y).
\end{equation}
The positive tail converges geometrically. On bounded intervals the
negative tail is dominated by a constant times $\sum_{j<0}(b\lambda)^j$.
For a coset $\omega+H_B$ write $m_{\omega+H_B}$ for its translate of
Haar probability. For $0<s\le1$ write $\mathcal S_s$ for the following
assertion. For every $0<h<s$, $T,A>0$ and every fixed upper bound on
the number of weight factors,
\begin{equation}\label{eq:anchor}
 \lim_{r\downarrow0}\ \sup_{\omega,v,c}
 m_{\omega+H_B}\!\left\{z:\int_{-T}^{T}
       \mathbf{1}_{\{|F_z(t)|\le Ar\}}\,dt>r^h\right\}=0.
\end{equation}
The supremum ranges over cosets, permitted rays and products of the
specified complexity at fixed multipliers and amplitude. For the
unweighted physical cosine alone, the same notation restricts
\eqref{eq:anchor} to $c_j=1$, $B=b$ and $v=1$. Recurring estimates
can increase that bound by one or more fixed units. Physical dilation
and reanchoring replace $v$ by $v(B/b)^n$ and $c_j$ by $c_{j+n}$.
They preserve the parameter families. A fixed phase translate absorbs
a change from $f$ to $f'$ apart from its fixed normalization.
More explicitly, write $F_z^{v,c}$ when the ray and weights are displayed.
For each integer $m\ge0$, reindexing both convergent tails gives
\begin{equation}\label{eq:fixed-dilation}
 F_z^{v,c}(b^m t)=\lambda^{-m}
 F_{S^{-m}z}^{\,v(B/b)^{-m},\, (c_{i-m})_{i\in\mathbb Z}}(t).
\end{equation}
Thus an anchored estimate on $[-1,1]$ with target radius $r$ gives
\eqref{eq:anchor} for any fixed $T,A>0$. Choose $m$ with $b^m\ge T$.
After the change of variable the occupation time is bounded by $b^m$
times the test on $[-1,1]$ at radius $\lambda^m Ar$. To obtain the
threshold $r^h$, apply that test with $h<h'<s$. The inequality
$b^m(\lambda^m Ar)^{h'}\le r^h$ holds for all sufficiently small $r$.
The shift transports coset Haar measure to coset Haar measure. For the
unweighted real cosine the transformed ray and weights are all one.

\begin{proposition}[Reciprocal transfer]\label{prop:transfer}
Suppose $\rho<\gamma=(b\lambda)^{-1}<1$. If $\mathcal S_s$ holds for the
dual system $(1/C,1/B,\gamma)$ at each fixed finite weight complexity,
then $\mathcal S_s$ holds for $(B,C,\lambda)$.
\end{proposition}
\begin{proof}
The past slope and its remainder for $0<|t|\le1$ are
\[
 H(z)=\sum_{l\ge1}\gamma^l c_{-l}p_{-l}(v,B)f'(z_{-l}),
 \qquad L_t(z)=F_z(t)/t-H(z).
\]
The new weight in $H$ depends on the fixed ray $v$ and
on $l$ alone. For some $h_0$ with
$\log\gamma/\log\rho<h_0<1$, uniformly in $t$ and $|h|\le1$,
\[
 |L_t(z+v_C(wh))-L_t(z)|\le C|h|^{h_0}.
\]
For the nonnegative indices use
$C|h|\sum_{j\ge0}(b\lambda\rho)^j$. For the negative-index Taylor
remainder the bound is
$C\sum_{l\ge1}\gamma^l b^{-l}|t|\min(1,\rho^{-l}|h|)$.
Choose $h_0$ sufficiently close to $\log\gamma/\log\rho$ that
$(\gamma/b)\rho^{-h_0}<1$. Splitting the series of $H$ gives
$|H(z+v_C(wh))-H(z)|\le C|h|^{\log\gamma/\log\rho}$.
At $h=\rho^m y$, $|y|\le2$, its increment divided by $\gamma^m$
is a dual increment, up to a fixed nonzero factor and an error
$O((\rho/\gamma)^m)$. The dual physical ray is $w(C/\rho)^{-m}$,
the environment is the shifted reversal $S^mRz$ with the constant
phase compensation for $f'$, and the weights at dual index $i$ are
$c_{-m-i}p_{-m-i}(v,B)$. Indeed $l=m+i$, so
$C^{-l}\rho^m=(C/\rho)^{-m}(1/C)^i$, and the unrepresented
indices $i\le-m$ contribute $O(\sum_{i\le-m}(\gamma/\rho)^i)$.

We now select the environments for the estimate. Take $0<\sigma<s$ and
$\kappa=1/K$ with $K\ge3$. Split $r^{1-\kappa}<|t|\le1$ into
bands $r^{(j+1)\kappa}<|t|\le r^{j\kappa}$, $0\le j\le K-2$,
and set $\Delta=r^{1-(j+1)\kappa}$. Take
$m_D=\lfloor\log\Delta/\log\gamma\rfloor$ and
$m_i=\lceil h_1m_{i+1}\rceil$ backwards, where
$h_1=\log\gamma/(h_0\log\rho)<1$. Fix $D$ with $h_1^D\le\kappa$.
The differences $d_i=m_{i+1}-m_i$ tend to infinity and
\[
 \rho^{h_0m_i}\le\gamma^{m_{i+1}},\quad
 (\rho/\gamma)^{m_i}\le\gamma^{d_i},\quad
 m_0\le\kappa m_D+O(1),\quad \Delta\asymp\gamma^{m_D}.
\]
For each of these finitely many band/stage choices test $\mathcal S_s$ on
the displayed dual environment, ray and weights at radius
$A\gamma^{d_i}$ with exponent $\sigma$ and $|y|\le2$.
Delete the failing original environments. The resulting set $G_r$
has complementary mass $o(1)$ uniformly. The choice of $G_r$ is independent of the particular
secant $t$ and of the transverse translation parameter.

Fix a starting environment $z_0$ and $t$ in the chosen band. For $h\in[0,1]$ set $z=z_0+v_C(wh)$ and
retain $h$ for which $z\in G_r$ and $|F_z(t)|\le r$.
Use nested dyadic intervals of lengths
$2^{-\lceil m_i\log_2(1/\rho)\rceil}\asymp\rho^{m_i}$.
Take a retained $h_*$ in an occupied parent. For every point $h'$
of any occupied child the preceding moduli and the two target
inequalities at $h_*$ and a retained point of that child imply
\[
 |H(z_0+v_C(wh'))-H(z_0+v_C(wh_*))|
 \le C(\gamma^{m_{i+1}}+r/|t|+\rho^{h_0m_i})
 \le C'\gamma^{m_{i+1}} .
\]
The good dual test at $z_0+v_C(wh_*)$, with
$h'=h_*+\rho^{m_i}y$, shows that the total length of those children
is at most $\rho^{m_i}\gamma^{\sigma d_i}$. Iterate to bound the
entire retained length by
$2^D\gamma^{\sigma(m_D-m_0)}\ll\Delta^{\sigma(1-\kappa)}$.
Average in $z_0$ by Haar invariance and integrate in the bands.
The retained mean occupation is
$O(r^{\sigma(1-\kappa)^2}+r^{1-\kappa})$. Markov's inequality,
then $\sigma\uparrow s,\ \,\kappa\downarrow0$, proves the claim on $[-1,1]$. A fixed physical rescaling gives each bounded test interval, and replacing the target radius by a fixed multiple preserves every strict exponent.

\end{proof}

\section{An initial exponent of one half on a physical ray}\label{sec:ray-half-seed}
Let $B,C$ denote the expanding and contracting multipliers of the
two embeddings under consideration. Set $L=\log|B|$, $M=-\log|C|$,
$\lambda=e^{-a}$, and assume
\[
          0<a<L,\qquad 0<a+M-L<L.                         \tag{H1}
\]
In particular (H1) holds when $M\le L$ and $0<L-a<M$.
Use the physical space $E=\mathbb R$ for a real $B$ and
$E=\mathbb C$ for a nonreal $B$. Write $v_B(y)_k=\operatorname{Re}(B^k y)$
and $w_C(u)_k=\operatorname{Re}(C^k u)$, with coordinates taken modulo one
when used as phases. The phase group has both these translation directions.
Indeed a Laurent character vanishing on the physical phase group satisfies
$p(B)=0$, and the conjugate embedding gives $p(C)=0$. Conversely, closure of
the entire ambient $v_B(E)$ has annihilator $\{p:p(B)=0\}$ also in the
complex case. The argument below requires invariance under the ambient
translations and does not require a single ray to be dense in that group.

Fix a unit physical ray $d$ and a unit transverse ray $e$ (take $1$ in a
real coordinate). Define
\[
 q_k=\operatorname{Re}((B/|B|)^k d),\quad
 p_k=\operatorname{Re}((C/|C|)^k e),\quad f(t)=\cos(2\pi t).
\]
For the real weights $c_k$ assume $|c_k|\le1$ and the following uniform
condition on positive indices. There are $D<\infty$, $\delta>0$ such that
every $D$ successive indices contain $k$ with
\[
                       |c_k q_k p_k|\ge\delta .              \tag{H2}
\]
For the uniform assertions the family of $(c_k,d,e)$ can be replaced
by its compact closure in the product topology times the two unit spheres.
Condition (H2) passes to this closure. On a phase coset take Haar $z$ and set
\[
 F_z(t)=\sum_{k\in\mathbb Z}\lambda^k c_k
    [f(z_k+v_B(td)_k)-f(z_k)].
\]
\begin{lemma}[Initial exponent]\label{lem:half}
Under (H1) and (H2), for each $0<s<1/2$, each fixed $K>0$, and $U$ uniform on $[0,T]$
independently of $z$, with $1\le T\le3$, one has
\[
 \mathbb P_{z,U}\{\nu_z([F_z(U)-Kr,F_z(U)+Kr])>r^s\}=o(1),
 \qquad \nu_z=(F_z)_*\operatorname{law}(U).                   \tag{H3}
\]
The limit is uniform over the indicated family, the phase coset and $T$.
For the permitted products of angular factors this gives $\mathcal S_{1/2}$.
\end{lemma}
\begin{proof}
We first estimate the probability in (H3), then pass to increments from the origin.

\par\smallskip\noindent\textbf{Transverse differences on the ambient space.}\enspace
Fix distinct real numbers $h_i,h_j$ in $[0,1]$. Let $D_z(y)$ be the
ambient difference of the two functions at $z+w_C(h_i e)$ and
$z+w_C(h_j e)$, with argument $y\in E$ in place of $td$. Its derivative
$G_z(y)=\partial_dD_z(y)$ exists and has an absolutely and uniformly
convergent plane-wave series. In fact the positive-frequency coefficient
at the linear functional $v_B(\cdot)_k$ has modulus
\[
 \pi\lambda^k |B|^k |c_k q_k|\,
 \left|\exp(2\pi i |C|^k p_k(h_i-h_j))-1\right|.
\]
For $k\ge0$ this is $O(Q^k)$, where $Q=\lambda|B||C|\in(0,1)$.
For all sufficiently large $k$ satisfying (H2) it is at least $cQ^k$,
with $c>0$ depending on the pair of translates. The negative time
coefficients are $O((\lambda |B|)^k)$ and their differentiated series
is bounded by $C\sum_{k<0}(\lambda |B|^2)^k$. Thus the derivative series
represents the derivative of the difference, by convergence on compact
sets of the original series and uniform convergence of the differentiated
series. Put $\eta=-\log Q/L\in(0,1)$. Splitting the positive series where
$|B|^k|y-y'|=1$ proves $|G_z(y)-G_z(y')|\le C|y-y'|^\eta$ globally.
All constants just used are independent of $z$ and of the allowed profiles
and ray directions.

The covector giving $v_B(\cdot)_k$ has Euclidean length $|B|^k$,
also when $E=\mathbb C$ is identified with $\mathbb R^2$. A fixed
Schwartz kernel on $E$, rotated to that covector and rescaled by $|B|^k$,
therefore isolates its coefficient. More explicitly, take a smooth
Fourier bump equal to one at the first unit vector, supported in a ball
about it of sufficiently small fixed radius to miss its negative and all
spheres of radii $|B|^j$, $j\ne0$. The inverse transform has integral
zero. At every center and every good $k$, the modulus of its convolution
with $G_z$ is at least $cQ^k$. Subtracting the central value, the
contribution outside distance $A|B|^{-k}$ is bounded by
$CQ^k\int_{|u|>A}|u|^\eta|\psi(u)|\,du$,
uniformly for rotated kernels (or use the radial envelope of their
absolute values). At fixed sufficiently large $A$ this is at most
$cQ^k/2$. Thus somewhere within this ball the derivative differs from
its central value by $c'Q^k$, and the H\"older estimate preserves half
that difference on a ball of radius $c''|B|^{-k}$. At a zero of $G_z$
these balls are holes of fixed relative volume at arbitrarily small
scales. Lebesgue differentiation gives $|\{y\in E:G_z(y)=0\}|=0$.
On every fixed bounded cube $J\subset E$ it follows that
\[
 \sup_{z,c,d,e}|\{y\in J: |G_z(y)|\le\kappa\}|\longrightarrow0
       \quad(\kappa\downarrow0).                              \tag{H4}
\]
A countersequence would have a convergent subsequence of phases, ray
directions and weights. The summable tails and continuity of finitely
many terms imply uniform convergence on $J$ to a member of the family.
All phase strings are allowed in (H4).

\par\smallskip\noindent\textbf{From area to the sampled line.}\enspace
The translation identity is $G_{z+v_B(y)}(td)=G_z(y+td)$, because the
constants subtracted on anchoring disappear on differentiation. Integrate
$y$ uniformly on a unit cube of $E$, and translate $z$ by $v_B(y)$
inside its coset Haar integral. For each $0\le t\le3$ this proves from
(H4), using a fixed cube containing all translates of the unit cube, that
\[
  \sup_{0\le t\le3}\mathbb P_z\{|G_z(td)|\le\kappa\}=o_\kappa(1).
                                                                    \tag{H5}
\]
The estimate is uniform in the coset, the weights and the directions.
It requires no assertion about restriction of an area-null set to an
individual line.

\par\smallskip\noindent\textbf{Pairs of heavy image cells.}\enspace
For a fixed integer $N$ choose $N$ distinct $h_i\in[0,1]$. Let $E_i(z)$
be the set of input points whose realized value for the phase
$z+w_C(h_i e)$ satisfies the heavy-ball test in (H3). The heavy output
points are contained in $O_K(r^{-s})$ intervals of length $r$, since the
enlargements of a grid of such intervals have bounded overlap and each
enlargement in question has mass exceeding $r^s$. Hence on
$E_i(z)\cap E_j(z)$, the scalar function $D_z(td)$ belongs to a union
of $O_K(r^{-2s})$ intervals of length $O(r)$. This union may depend on
$z,i,j$. Divide $[0,T]$ into a fixed number $O_\kappa(1)$ of intervals
of diameter small enough that the oscillation of $G_z(td)$ on each
interval is at most $\kappa$. On any interval containing a point where
$|G_z(td)|>2\kappa$, the derivative of the scalar function has fixed
sign and absolute value at least $\kappa$. The intersection with the
two heavy sets there has length $O_\kappa(r^{1-2s})$. The remaining
intervals lie in $|G_z(td)|\le2\kappa$, of mean length $o_\kappa(1)$
by (H5). First let $r\downarrow0$, then $\kappa\downarrow0$, to see that
the mean overlap of each pair tends to zero uniformly. Since at each
point $\sum_i{\bf1}_{E_i}\le1+\sum_{i<j}{\bf1}_{E_i\cap E_j}$ and
Haar invariance makes all $N$ summands on the left have the same mean,
the limsup of the probability in (H3) is at most $1/N$. Let $N\to\infty$.
For the anchored version, take $T=3$ and $1\le t\le2$. The centered
mass at $F_z(t)$ dominates one third the anchored mass on $[-1,1]$
for $z+v_B(td)$; that reanchored state has coset Haar law. A strict
exponent loss absorbs the factor three.

\par\smallskip\noindent\textbf{Products of angular factors.}\enspace
Here at most one of the two embeddings is nonreal. Write its angle as
$\theta$. Apart from real sign factors all the factors tested in
$c_k q_k p_k$ are of the form $\cos(\pm k\theta+\phi)$.
For irrational $\theta/(2\pi)$, a product of a fixed number $D'$ of
these factors is a nonzero trigonometric polynomial for every choice
of its phases. The integral of its squared modulus therefore has a
positive minimum over the compact phase torus. Uniform convergence of
the orbit averages of this squared trigonometric polynomial follows
directly by summing the nonconstant Fourier modes as geometric
progressions, uniformly in the phases and the starting index.
This proves (H2) with constants depending on $D'$ and $\theta$.
For rational rotation choose the complex ray $e^{i\phi_0}$ with
$\cos(k\theta+\phi_0)\ne0$ on every residue. This is possible by
avoiding finitely many directions. Its conjugated and rotated
directions also have a positive lower bound on their projection
factors. A product of a fixed number of these factors then meets
(H2). More generally, for any normalized periodic profile with
$\max_k|c_k|=1$ and common period the order of the rotation
(enlarged by a factor two for real signs if needed), (H2) holds
uniformly by using a residue attaining that maximum. This last
form includes compact limits of normalized periodic products.
\end{proof}
\section{The weighted critical estimate on a fixed ray}
\label{sec:weighted-critical-ray}
Write $b=|B|=e^L>1$, $\rho=|C|=e^{-M}<1$ and
$\lambda=e^{-a}$. In this section $a+M=L$, equivalently
$\lambda b\rho=1$. Thus $0<M<L$ and $\gamma=(b\lambda)^{-1}=\rho$.
The physical frequency $B$ and transverse frequency $C$ can be real or
complex. The Haar group has the following two properties: a finitely
supported character with integer coefficients $n_j$ is trivial precisely
if $\sum_j n_j B^j=0$, and the translations with coordinates
$\operatorname{Re}(C^j h)$ preserve Haar measure on each coset. In a
complex physical space the first property refers to the group of the
\emph{full} physical space before choosing a ray. For physical complex
arguments we use the convention $\operatorname{Re}(B^j x)$.

Let $v,w$ be fixed unit physical and transverse directions and put
\[
 d_j=\operatorname{Re}((B/b)^j v),\qquad
 e_j=\operatorname{Re}((C/\rho)^j w),\qquad
 k_j=c_jd_je_j .
\]
Here $c_j$ are the real weights of the tested model, with $|c_j|\le1$.
We allow arbitrary additive forcing phases by working on translated Haar
cosets. The estimate below is uniform in any family of weights and rays
for which, for some $\eta>0$ and $m_0<\infty$,
\begin{equation}
 \inf_{i\in\mathbb Z}\sum_{j=i}^{i+m-1} k_j^2\ge\eta m
 \qquad(m\ge m_0). \tag{RC1}
\end{equation}
For example, with at most one nonreal direction, finite products of
cosines are trigonometric polynomials evaluated along one rotation.
For an irrational rotation, any nonzero polynomial of this type has
positive square mean; multiplying by any additional cosine preserves
this fact. The geometric series for each nonconstant Fourier mode gives
uniform convergence of the block square means, independently of their
origins. On compact normalized families of fixed degree this is uniform
by compactness. For a finite cyclic rotation one instead takes the mean
over its actual finite orbit. One can choose the physical or transverse
ray in a complex place to avoid the finitely many perpendicular
directions, making the additional factor nonzero at every point of the
orbit. Simultaneous translations of the whole sequence $k_j$ preserve
(RC1). For a fixed arbitrary ray the stated estimate requires (RC1);
an identically zero product on a finite orbit would invalidate its
variance bound.

Write $V_z(tv)=F_z(t)$ for the weighted increment along the physical
ray $v$, with $f(t)=\cos(2\pi t)$ as in \eqref{eq:increment}.

\begin{proposition}[Critical estimate]\label{prop:critical}
Assume $a+M=L$ and (RC1). For every fixed $A,K>0$ and $0<s<1$, 
\begin{equation}
 m\left\{z:\int_{-A}^A {\bf1}_{|V_z(tv)|\le Kr}\,dt>r^s\right\}
 \longrightarrow0. \tag{RC2}
\end{equation}
Here $m$ is Haar probability on the chosen coset. For fixed $B,C$ the
convergence is uniform over cosets and over weights and rays with common
constants $\eta,m_0$ in (RC1). In particular $\mathcal S_1$ holds for the permitted products.
\end{proposition}
\begin{proof}
We take $A=1$. For fixed $A$ the same calculations take finitely
many additional negative values of the physical scale index.

\par\smallskip\noindent\textbf{Concentration of the drift.}\enspace
Define real sums $S_0=0$ on both sides of zero by
$S_{j+1}-S_j=k_j f''(z_j)$. We claim that for each fixed $D>0$,
\begin{equation}
 \sup_i\sup_{t\in\mathbb R}m\{|S_{i+m'}-S_i-t|\le D\}\to0
       \quad(m'\to\infty). \tag{RC3}
\end{equation}
For completeness, normalize $f''$ to a cosine. The variance on the
block is $\sigma^2=\frac12\sum k_j^2\asymp m'$, since distinct powers
of $B$ have different absolute values. Only vanishing signed sums of
powers of $B$ survive in the Haar integral of a product of characters.
For each fixed $q$ all nonzero sums of at most $q$ signed nonnegative
powers have absolute value at least some $c_{q,B}>0$. If this failed,
pass to a countersequence with fixed signs and sorted exponents, divide
by the largest complex power, and pass to a subsequence on which all
offsets from that exponent are fixed or diverge. The fixed-offset group
sums to zero exactly. Delete it and apply induction on $q$ to the
remaining nonzero sum. Bounded largest exponents give only finitely many
possible sums. It follows that any such zero relation separates across
every gap exceeding a fixed $J_q$. If $m$ is the smallest exponent of the upper group, its nonzero sum has modulus at least $c_{q,B}b^m$, by factoring out $B^m$. The lower group has modulus at most $q b^{m-J_q}$. Choose $J_q$ with $qb^{-J_q}<c_{q,B}$; a zero relation then forces each group to sum to zero. A cluster has at least two occurrences; clusters of size two
are equal times with opposite signs. For the $2l$-th moment there are
$O_{l,B}((m')^{l-1})$ configurations having a cluster of size bigger
than two. Pairing all the others gives
$(2l-1)!!\,\sigma^{2l}+O_{l,B}((m')^{l-1})$, since choices of pairs
with nearby indices also have this error order. For order $2l+1$ the
moment is $O_{l,B}((m')^l)$. The bounds use only $|k_j|\le1$ and hold
for every coset. The normalized laws therefore converge to the standard
normal along every sequence of allowed blocks and weights, proving
(RC3) by tightness and continuity of its distribution function.
Restoring the fixed factor in $f''$ has no effect. In particular,
for fixed $H,D$ the following quantity tends uniformly to zero in
probability as $n\to\infty$:
\[
 n^{-1}\sup_t\#\{i\in[-Hn,Hn]\cap\mathbb Z:|S_i-t|\le D\}.
\]
Indeed its square is bounded by $n^{-2}$ times the number of pairs with
$|S_i-S_j|\le2D$. Split into small and large index gaps and apply (RC3)
to see that the expected number of such pairs is $o(n^2)$.

\par\smallskip\noindent\textbf{Critical linearization.}\enspace
Write $W(h)_j=\operatorname{Re}(C^j w)h$. Given $0<|t|\le1$, choose
$k\ge0$ by $b^{-k-1}<|t|\le b^{-k}$. For $N\ge0$ and $|h|\le D\rho^N$
one has, uniformly,
\begin{equation}
 \frac{V_{z+W(h)}(tv)-V_z(tv)}{t}
       =h(S_{k+1}-S_{-N})+O_D(\rho^N). \tag{RC4}
\end{equation}
For $j<-N$ use the physical first difference bound $C b^j |t|$.
For $j>k$ use the transverse bound $C\rho^j |h|$ and
$(\lambda\rho)^k/|t|=O_b(1)$. For $-N\le j\le k$ integrate the mixed
difference twice. Its weighted quotient equals $h k_j f''(z_j)$ plus
an error of size at most $C|h|(b^j|t|+\rho^j|h|)$, since
$\lambda b\rho=1$ and $|d_j|,|e_j|,|c_j|\le1$. Summing gives (RC4).

\par\smallskip\noindent\textbf{Retaining environments before integrating in the secant.}\enspace
Put $n=\log(1/r)$, fix $0<\delta<1/4$, and for $r^{1-\delta}<|t|\le1$
put $T=Kr/|t|$. Choose a large fixed $q$, let $d=2^{-q}$, and
subdivide the unit transverse interval into nested intervals of length
$d^l$. Let $N_l$ satisfy $\rho^{N_l}\asymp d^l$ by nearest-integer
rounding, and take $J=\lfloor\log(1/T)/\log(1/d)\rfloor$. For large
$q$ the $N_l$ are distinct. All indices $-N_l,k+1$ belong to
$[-Hn,Hn]$ with $H$ fixed. For $D_q$ a sufficiently large constant
times $d^{-1}$ choose a set $G_r$ of measure $1-o(1)$, uniformly in
the model, on which at most $\epsilon n$ of those indices can lie in
\emph{any} interval of radius $D_q$ in the value of $S_i$. This set
is chosen independently of $t$, by the consequence of (RC3).

Fix $z_0$ in the coset and consider the $h_0\in[0,1]$ for which
$z=z_0+W(h_0)\in G_r$ and $|V_z(tv)|\le Kr$. At such a point mark
level $l<J$ bad if $|S_{k+1}-S_{-N_l}|\le D_q$. There are at most
$\epsilon n$ bad marks. Fix one pattern of marks. In an occupied
parent at a good level apply (RC4) at one selected reference point to
every selected point in the same parent. Their distances from that
reference are bounded by $C(T+d^l)/D_q=O(d^{l+1})$. Thus at most a
fixed $C_*$ children are occupied. This constant is independent of $q$:
rounding $N_l$ makes $d^l/\rho^{N_l}$ bounded in terms of $\rho$ alone,
so (RC4) is used with fixed $D$ throughout. At a bad level allow all $d^{-1}$.
With at most $2^J$ patterns, the selected length is bounded by
$(2C_*d)^Jd^{-\epsilon n}$. Since $\log(1/T)\ge\delta n-O_K(1)$,
first choosing $q$ large, then $\epsilon$ small, bounds this length by
$C_{q,K} T^{1-\delta}$. Average over $z_0$ and use Haar invariance
under $W(h_0)$. Integrating in $t$, with the trivial bound on
$|t|\le r^{1-\delta}$, yields
\[
 \int_{G_r}\int_{-1}^1 {\bf1}_{|V_z(tv)|\le Kr}\,dt\,dm(z)
      \le C_{q,K,\delta}r^{1-\delta},
\]
since $\int_0^1t^{-(1-\delta)}dt<\infty$. Markov's inequality and
$m(G_r^c)=o(1)$ prove (RC2) for every $s<1$.


\end{proof}

\section{Equal rates and the strip estimate}\label{sec:equal}

Assume $L=M$, so $b\rho=1$, and write
$\alpha=a/L$, $\beta=1-\alpha$, $\gamma=(b\lambda)^{-1}=b^{-\beta}$.
\begin{proposition}[Improvement at equal rates]\label{prop:strip}
Suppose the original system has $\mathcal S_s$ and the reversed
system at amplitude $\gamma$ has $\mathcal S_t$. Then the original
system has
\[
 \mathcal S_{\beta s+(1-\beta)\min(1,\beta s+t)}.               \tag{B1}
\]
The assertions apply to the parameter families at each fixed finite
weight complexity.
\end{proposition}

\begin{proof}
Take $0<s'<s$, $0<t'<t$ and an integer $J\ge2$.
At $r=b^{-m}$ put $\Delta=\gamma^m=r^\beta$, $d=r/\Delta$,
and use the meshes $l=d^{j/J}$, $0\le j\le J$.
Independently sample Haar-coset $z$, uniform stable parameter
$h\in[0,1]$ and uniform physical input $x\in[0,3]$. Divide the
$h$ interval into $\lceil1/r\rceil$ equal cells with centers $h_0$.
Set $Y(x)=F_{z+v_C(wh_0)}(x)$ and
\[
 P(h)=\sum_{k\ge1}\gamma^k c_{-k} p_{-k}(v,B)
             f'(z_{-k}+\operatorname{Re}(w C^{-k})h).
\]
Uniformly for $x\in[0,3]$ and $h$ in its cell,
\[
 F_{z+v_C(wh)}(x)=Y(x)+x(P(h)-P(h_0))+O(r).                    \tag{B2}
\]
The positive-index variation is at most
$Cr\sum_{k\ge0}\lambda^k\rho^k$; for the negative indices subtract
the part linear in $x$ and differentiate the remainder in the stable
phase, for a bound $Cr\sum_{k\ge1}\lambda^{-k}b^{-2k}\rho^{-k}$.
Both series converge. Splitting the slope series at $m$ gives
$|P(h)-P(h_0)|\ll\Delta$.

Condition on $z$ and the stable cell and write
$\theta=(P(h)-P(h_0))/\Delta$. Discard stable parameters of total
mass $o(1)$ on average so that the resulting conditional direction
subprobability $\sigma$ satisfies
\[
 \sigma(I)\ll l^{t'} \quad (|I|\le C l)                         \tag{B3}
\]
at every listed mesh. Indeed, for a sampled $h_*$, change variables
$h=h_*+b^{-m}y$, $|y|\le1$, in the centered conditional test. The
normalized density is bounded and, by the exact reversed-slope
computation above, the normalized slope increment is a dual increment
at $S^m R(z+v_C(wh_*))$ (up to fixed phase adjustment and multiplicative
factor), with error $O((\rho/\gamma)^m)=O(d)$. The transformed rays
and weights are deterministic for this test, and the environment is
translated Haar for each fixed $h_*$. At the $J$ shrinking meshes
$l\ge d$ the dual assertion $\mathcal S_t$ applies using slack above
$t'$. The mesh $1$ is trivial. Remove the union of failing centered
tests. A target meeting the retained directions is contained in a
doubled test at a retained point. The resulting bound averages the
deleted conditional stable probability over $z$ and the cell choices.

Independently of the retained parameter within its cell, discard an
averaged $o(1)$ of the physical inputs for $Y$. For each range strip
$[i\Delta,(i+1)\Delta)$ let $\mu_i$ be the retained uniform input
measure pushed to coordinates $(x,(Y(x)-i\Delta)/\Delta)$ on that
strip and write $M_i=\|\mu_i\|$. At the listed meshes,
\[
 \sup_p\mu_i(B(p,l))\ll\Delta^{s'}l^{D},\qquad
 D=1+\beta s',\qquad M_i\ll\Delta^{s'},\quad\sum_iM_i\le1.     \tag{B4}
\]
In fact for sampled input $x_*$ and $b^{-n}\asymp l$, recenter at
$x_*$, rescale by $b^{-n}$, and apply the original $\mathcal S_s$ to
the requirement $|x-x_*|\le C l$,
$|Y(x)-Y(x_*)|\le C\Delta l$. The normalized input is bounded, the
normalized target radius is comparable to $\Delta l^\beta$, the new
physical ray is $v(B/b)^n$, and the new weights are $c_{j+n}$.
The Haar marginal after the two translations and shift holds at
each fixed $h_0,x_*$. All the target radii tend to zero. Removing
the finite bad tests gives their bounds at any retained center,
and hence arbitrary centers via a doubled test. The mesh $1$ gives
the stated strip mass since the normalized strips have bounded
support. With the environment and the stable cell fixed, the restriction
defining $\sigma$ depends only on $h$. Those defining the $\mu_i$ depend
only on $x$ and the cell center $h_0$. Thus the resulting measures may
be integrated against one another before averaging over the cells.

We give the frequency estimate since the factor $\Delta^{s'}$ must
be retained at all frequencies. For each $h$, two retained outputs
within $Cr$ give by (B2) two strip points with
$|i-i'+\pi_\theta(p)-\pi_\theta(p')|\le C' d$, where
$\pi_\theta(x,y)=y+\theta x$. Only $|i-i'|\le C''$ contribute.
Use a positive Schwartz majorant of this interval indicator with
Fourier support in $[-C/d,C/d]$ and Fourier norm at most $C d$
pointwise. One obtains it by rescaling a squared modulus of an
inverse Fourier transform of a nonnegative smooth bump. Taking
absolute values and using $2yz\le y^2+z^2$ bounds the conditional
averaged collision probability by
\[
 C d\sum_i\int\int_{|\tau|\le C/d}
         |\widehat\mu_i(\tau\theta,\tau)|^2\,d\tau\,d\sigma(\theta).
\]
The range $|\tau|\le C$ contributes $O(d\Delta^{s'})$. On every
larger band $C d^{1/J}/l\le|\tau|\le C/l$, the image of
$d\sigma(\theta)d\tau$ under $(\theta,\tau)\mapsto(\tau\theta,\tau)$
has unit-square masses $O(d^{-1/J}l^{t'})$ by (B3), because the
allowed angular interval at each $\tau$ is covered by $O(d^{-1/J})$
intervals of length $l$. Its support lies in $B(0,C'/l)$.
Uniformly in $y\in\mathbb R^2$, the Gaussian integral majorization
and Fourier inversion give
\[
 \int_{B(y,C'/l)} |\widehat\mu_i|^2
  \ll M_i\Delta^{s'}l^{D-2}.
\]
The physical kernel is bounded by $Cl^{-2}\exp(-c|p-p'|^2/l^2)$,
and a ring of radius $kl$ is covered by $O((k+1)^2)$ balls of
radius $l$. To use this disk estimate against the band law, majorize
$|\widehat\mu_i|^2$ by $C_N(1+|\cdot|)^{-N}*|\widehat\mu_i|^2$:
take a smooth cutoff identically one on the common compact physical
support, convolve on the frequency side and use Cauchy--Schwarz.
Convolving the band law with this kernel produces a density bounded
by $C_N d^{-1/J}l^{t'}(1+\operatorname{dist}
(\cdot,B(0,C'/l)))^{-N+3}$. Sum the uniform disk estimate over
boxes of side $1/l$, taking $N$ large. The bound for an individual strip
still contains $M_i\Delta^{s'}$. Since $\sum_iM_i\le1$, after the initial
factor $d$ the cost of a band is at most
$C d^{1-1/J}\Delta^{s'} l^{\,\beta s'+t'-1}$.
For a nonnegative exponent of $l$ use $l\le1$; for a negative exponent
use $l\ge d$. Summing the $J$ bands and the low frequencies gives
\[
 C_J\Delta^{s'}d^{\min(1,\beta s'+t')-1/J}.                      \tag{B5}
\]

To identify the collision estimated in (B5), fix $z$ and a transverse
cell $J_0$. Let $U\subset J_0$ and $Q\subset[0,3]$ be the retained
transverse and physical inputs above. Write $m_{J_0}$ and $\ell$ for
normalized Lebesgue measure on these two full intervals. For
$h\in J_0$ put $\zeta=z+v_C(wh)$ and set
\[
 \nu_\zeta=(F_\zeta)_*\ell,\qquad
 \eta_{z,J_0,h}={\bf1}_U(h)(F_\zeta)_*(\ell|_Q)\le\nu_\zeta .
\]
With $\mathcal C_R(\eta)=(\eta\otimes\eta)\{(y,y'):|y-y'|\le R\}$,
the conditional collision has the exact representation
\[
 \int_{J_0}\mathcal C_R(\eta_{z,J_0,h})\,dm_{J_0}(h)
 =\int_U\int_{Q^2}
 {\bf1}_{\{|F_\zeta(x)-F_\zeta(x')|\le R\}}\,
 d\ell(x)\,d\ell(x')\,dm_{J_0}(h).
\]
Both ordinates on the right use the same $h$. At $R=Cr$, (B2) and
the strip estimate above bound the right-hand side by (B5).
On averaging over $z,J_0,h$, the omitted mass is $o(1)$ by the two
deletion estimates. Lemma~\ref{lem:full-ball} below applies separately
to each pair $(\nu_\zeta,\eta_{z,J_0,h})$, including different
representations of the same $\zeta$. Take $\tau=r^u$ with $u$ strictly
below the exponent in (B5) expressed in $r$. At fixed $M$, the
averaged collision term divided by $\tau$ tends to zero. Let $M\to\infty$,
then reanchor
sampled $x\in[1,2]$; its centered test on $[0,3]$ dominates the
anchored occupation on $[-1,1]$ in the Haar-coset environment
$z+v_C(wh)+v_B(vx)$. Extend to each fixed $T$ by a fixed physical
rescaling, and to arbitrary radii by geometric interpolation. This
proves (B1) by $s'\uparrow s$, $t'\uparrow t$, $J\to\infty$.

\end{proof}

\begin{lemma}[Full ball mass from retained collisions]\label{lem:full-ball}
Let $\nu$ be a probability measure on $\mathbb R$ and let $0\le\eta\le\nu$.
Write $\mathcal C_R(\eta)=(\eta\otimes\eta)\{|y-y'|\le R\}$.
For $R,\tau>0$ and $M>1$,
\[
 \nu\{y:\nu([y-R,y+R])>\tau\}
 \le 1-\|\eta\|+\frac{3}{M}+\frac{M}{\tau}\mathcal C_R(\eta).
\]
\end{lemma}
\begin{proof}
Partition the line into half-open cells $I$ of length $R$. If
$\nu(3I)>M\eta(I)$, summing $\eta(I)$ over these cells costs at most
$3/M$, since the concentric triples have overlap three. In any other
cell containing a point of the displayed event,
$\tau<\nu(3I)\le M\eta(I)$. Its retained mass is at most
$(M/\tau)\eta(I)^2$. Pairs from the same cell have distance at most
$R$, so $\sum_I\eta(I)^2\le\mathcal C_R(\eta)$. The contribution of
$\nu-\eta$ is at most $1-\|\eta\|$.
\end{proof}

Both systems satisfy $\mathcal S_{1/2}$ by Lemma~\ref{lem:half},
since in each orientation $0<a+M-L=a<L$. If every finite
complexity admits a common exponent $u$ in both directions, (B1)
improves it to the smaller of
$p u+(1-p)\min(1,(1+p)u)$ for $p=\alpha,\beta$.
Both expressions are strictly greater than $u$ for $0<u<1$.
Iteration from $1/2$ therefore reaches every exponent below one.
Each probability estimate involves a fixed finite weight complexity.

\section{Transverse differentiation and the amplitude recursion}\label{sec:descent}
\begin{proposition}[Transverse differentiation]\label{prop:descent}
Suppose $a+M<L$. A uniform
$\mathcal S_s$ for $(B,C,e^{-(a+M)})$ implies
\[
 \mathcal S_{\min(1,(a+M)s/a)}\quad\text{for }(B,C,e^{-a}).       \tag{D}\label{eq:descent}
\]
\end{proposition}
\begin{proof}
Write $g=e^{-a}$ and $h=e^{-(a+M)}$. Differentiating $F_{z+v_C(wu)}(x)$ in $u$ produces, up to a fixed nonzero constant, the new function $G$ with amplitude $h$, weights $c_jp_j(w,C)$ and a constant phase offset. For $j\ge0$ the series of derivatives is dominated by $C h^j$; for $j<0$ it is dominated by $C(h|B|)^j$ on bounded physical intervals. Uniform convergence gives a uniformly continuous derivative in $u$ on those intervals, uniformly in phase and permitted weights/rays, by truncating the two tails. The new inputs stay in the parameter families declared above.

For completeness we give the information argument for (D). On $[0,T]$, $1\le T\le3$, use normalized Lebesgue measure and coset Haar phase. An anchored bound at exponent $s$ for $G$ implies the sampled bound at every strict exponent $s'<s$: recenter at an independent uniform $x$ and use the anchored interval $[-3,3]$. Fix $0<t<\min(1,(a+M)s/a)$ and a large fixed block size $n$. Divide the interval into $N=\lfloor T|B|^n\rfloor$ equal pieces. Rescaling a piece to $[0,T']$, $T'=T|B|^n/N\in[1,2]$, scales increments of $F,G$ by $g^n,h^n$ respectively, shifts the weights by $n$, rotates the ray, and maps the phase to a translated Haar environment. The oscillation of $G$ on any such piece is $\le C_0 h^n$. Join pieces when their values of $G$ at the left endpoints are within $(2C_0+1)h^n$. Count the piece itself among its neighbors. Mark a piece bad if its degree exceeds $Ng^{nt}$. For an input in that piece all neighboring pieces map by $G$ inside its ball of radius $(4C_0+1)h^n$. The sampled bound, with $s'$ such that $t a/(a+M)<s'<s$, gives probability $e_n=o(1)$ of a bad piece, uniformly in the current deterministic parameters. Repeat this construction in the rescaled child with its updated parameters; the test at each node uses the derivative $G$ for the newly started $F$. Only the magnitude and family of that derivative are needed. Fix the same admissible transverse ray at every node of the orientation under discussion; the child suffixes then use precisely the same rule.

To see how the transverse direction enters a child, denote the derivative
under $z\mapsto z+v_C(wu)$ at $u=0$ by $D_w$. For a child with left endpoint
$x_0$ put $\widetilde z=S^n(z+v_B(vx_0))$. Let $\widetilde F$ have
the same amplitude $g$, physical ray $v(B/b)^n$ and weights $c_{i+n}$.
Reindexing the two-sided increment gives
\[
 F_z(x_0+b^{-n}y)-F_z(x_0)=g^n\widetilde F_{\widetilde z}(y).
\]
Under transverse translation $\widetilde z$ changes by $v_C(C^n w u)$.
Differentiating both sides therefore gives
\[
 D_w F_z(x_0+b^{-n}y)-D_w F_z(x_0)
    = h^n D_{(C/\rho)^n w}\widetilde F_{\widetilde z}(y).
\]
This identity bounds the oscillation used in the parent's neighbor test,
uniformly in the rotated direction. The child's own neighbor test uses
$D_w\widetilde F$ by the recursive definition. The diagonal energy below
rescales $F$ and the statuses generated by those child tests.

For a length-$k$ status word $P$ retain the unnormalized input on the successive labels prescribed by $P$. Write $E_P(r)$ for the supremum over allowed parameters (with the present fixed complexity bound for $F$), cosets and lengths of its mean self-energy under $K_r(y)=r^{-1}(1-|y|/r)_+$. If $P'$ is its suffix then
\[
 E_P(r)\le C_n+A E_{P'}(r/g^n),\qquad
 A=g^{-n(1-t)}\ \text{for a good first label},\quad
 A=g^{-n}\ \text{for a bad one}.
\]
Indeed, for a pair of non-neighboring children the derivative of the difference of the two realized ordinates with respect to transverse translation has modulus at least a fixed constant times $h^n$. Average over an extra Haar-preserving translation $u\in[0,1]$, discard all path restrictions there, and keep this necessary derivative bound at the translated environment. Uniform continuity of that derivative gives a finite partition of $[0,1]$ on each of whose intervals meeting the bound the derivative has constant sign and a uniform positive lower modulus. Change variables on it to bound the integral of the kernel by $C_n$, uniformly in the input pair. For neighbors use that $K_r$ is an autocorrelation and hence defines a positive semidefinite bilinear form. Bound each mixed energy by half the sum of the two diagonal energies. Its multiplicity in the good case is at most $Ng^{nt}$ since both labels are good; the unrestricted multiplicity is $N$. The child masses carry factors $N^{-2}$. A diagonal rescaling costs $g^{-n}$ in the kernel. Discard the first label condition in each diagonal and average in the phase for that deterministic child, whose transformed phase is translated Haar. This gives the displayed recursion. The rescaled model for $F$ acquires a shifted existing weight, with the same complexity; the extra weight for $G$ enters only its neighbor test.

At scale $r_k=g^{nk}$ retain all words with at most $\epsilon k$ bad labels. The discarded mass on average is at most $e_n/\epsilon$: before imposing any labels, the expected bad probability at each depth is at most $e_n$ by averaging the deterministic physical cells and their shifted Haar states. There are at most $\exp(kH(\epsilon))$ retained words for $0<\epsilon<1/2$, with binary entropy $H$. Iterating the recursion and applying the triangle inequality to the kernel seminorm of their sum bounds the retained collision probability at distance $Cr_k$ by
\[
 C_{n,C}(k+1)r_k^{\,t(1-\epsilon)-2H(\epsilon)/(na)} .
\]
Use the kernel at $2Cr_k$, where a close pair contributes at least $1/(4Cr_k)$. Apply Lemma~\ref{lem:full-ball} at $R=Cr_k$ to the full law and
the submeasure retained at this scale. For any exponent strictly below
the displayed collision exponent, take $\tau$ equal to $r_k$ to that
power, let $k\to\infty$ at fixed $n,\epsilon,M$, then let $M\to\infty$. Fixed $n$ allows interpolation between the $r_k$ with a fixed dilation. Choose $n\to\infty$, $\epsilon\downarrow0$ with $e_n/\epsilon\to0$, and reanchor on $[1,2]\subset[0,3]$ to obtain (D). Fixed larger anchored domains follow by dilating the physical input. All bounds used here are for fixed finite weight complexity.

\end{proof}

\subsection{The recursion at unequal rates}\label{sec:closure}
\begin{proposition}[A contracting conjugate]\label{prop:two}
Let $B,C$ be as in Section~\ref{sec:embeddings}, with
$|C|<1<|B|$, and let $|B|^{-1}<\lambda<1$.
Then $\mathcal S_1$ holds on the permitted families of rays and weights.
\end{proposition}
\begin{proof}
Suppose first that $L>M$, and write $\lambda=e^{-a}$.
At $a=L-M$ Proposition~\ref{prop:critical} gives $\mathcal S_1$.
If $0<a<L-M$, Proposition~\ref{prop:descent} uses the estimate at $a+M$.
If $L-M<a<L$, apply Proposition~\ref{prop:transfer} twice.
The intermediate log rates are $(M,L)$ at amplitude exponent $L-a$.
The second transfer returns to rates $(L,M)$ at exponent $a+M-L$.
The two hypotheses are $a+M>L$ and $(L-a)+L>M$.
The latter follows from $a<L$ and $L>M$.
Both transfers preserve the available information exponent.

Up to a visit to $a_n=L-M$, the sequence of required parameters is
\[
 a_n=(a+nM)\bmod L.
\]
In the first $n$ steps the number of indices with $a_i+M>L$ is
$\lfloor(a+nM)/L\rfloor$ in the absence of a critical visit.
Since $0<M/L<1$, there are infinitely many steps of each type.
At every index with $a_i+M>L$ we have
$0<a_i+M-L<M<L$. Thus Lemma~\ref{lem:half} gives
$\mathcal S_{1/2}$ there. At every earlier index with $a_i+M<L$,
\eqref{eq:descent} increases the exponent in the backward deduction
by a factor $(a_i+M)/a_i\ge1+M/L$, up to the cap of one.
Choose an index of the first type with at least $k$ earlier indices
of the second type, where $(1+M/L)^k/2\ge1$.
The backward deduction gives $\mathcal S_1$ at $a_0$.
If the orbit has a critical visit, start there with exponent one.
There is also a bound on the number of required parameter steps, uniform
in $a$ for fixed $L>M$. Set $r_0=M/L$,
$k_* = \lceil\log 2/\log(1+r_0)\rceil$ and
$N=\lceil(k_*+1)/(1-r_0)\rceil$. Unless the critical value is visited,
the first $N$ steps contain at least $N(1-r_0)-1\ge k_*$
indices with $a_i+M<L$, by the floor count above. The next
$\lceil1/r_0\rceil$ steps contain an index with $a_i+M>L$, since
adding $\lceil1/r_0\rceil r_0\ge1$ increases the floor by at least one.
At that index the initial exponent $1/2$ suffices to start the backward
deduction. An earlier critical visit starts it with exponent one.
Thus $N+\lceil1/r_0\rceil$ parameter steps suffice, each involving
only finitely many weight factors.

For $M>L$, a single reciprocal transfer applies since $a+M>L$.
In the reversed system $0<L-a<M$ and the larger log rate is first,
so the preceding case applies. For $L=M$, Proposition~\ref{prop:strip}
and its iteration in Section~\ref{sec:equal} give the conclusion.
\end{proof}

\appendix
\section{From anchored information to the graph}\label{app:graph}
\begin{proposition}[Graph comparison]\label{prop:graph}
Fix a real $b>1$, $1/b<\lambda<1$, and $0<s\le1$.
Suppose the unweighted cosine increment satisfies \eqref{eq:anchor}
uniformly on cosets of $H_b$. For every prescribed phase sequence and
every interval of positive length the Hausdorff dimension of the graph
is at least $1+s(1-\alpha)$. At $s=1$ the graph and its natural measure
satisfy the conclusions of Theorem~\ref{thm:main}.
\end{proposition}
\begin{proof}
For a prescribed one-sided sequence of phases, proceed as follows. Extend them in any way to all integer indices. The compact phase group $H$ is the closure of $t\mapsto(b^it\bmod1)_{i\in\mathbb Z}$ in $\T^{\mathbb Z}$, as in the two-embedding argument. Indeed its annihilator consists exactly of the Laurent polynomials vanishing at $b$. Write
\[
 \theta_n(x)=(b^{n+i}x+\varphi_{n+i})_{i\in\mathbb Z},
 \quad V(\theta,u)=\sum_{i\in\mathbb Z}\lambda^i
 [\cos(2\pi(\theta_i+b^iu))-\cos(2\pi\theta_i)].
\]
On $|u|\le1$ this is continuous by uniform convergence of both tails. It differs from $\lambda^{-n}(W(x+b^{-n}u)-W(x))$ by $O((b\lambda)^{-n})$. Splitting at $b^i|u|=1$ shows that $|V(\theta,u)|\le C|u|^\alpha$ and $|W(x+u)-W(x)|\le C|u|^\alpha$, uniformly in the phases.

Fix a continuous cutoff $0\le\psi\le1$ equal to one on $[-3,3]$, vanishing outside $[-4,4]$, and let
\[
 B_q(\theta)=2\int_{-1}^1\psi(V(\theta,u)/\lambda^q)\,du,
 \qquad f_q(\theta)=-\log\min(1,B_q(\theta)).
\]
The H\"older bound near $u=0$ gives a positive lower bound on $B_q$ for each fixed $q$, so $f_q$ is continuous and bounded. The definition of $\mathcal S_s$ implies, by restricting the negative logarithm to the good set,
\[
 \liminf_{q\to\infty}q^{-1}\inf_{\omega\in\T^{\mathbb Z}}
 \int_{H} f_q(z+\omega)\,dm_{H}(z)\ge s\log(1/\lambda).\tag{G1}
\]
Work on a translated unit interval with $X$ uniform there and $Y=W(X)$. Denote by $X_k$ the index of the horizontal dyadic interval of length $2^{-k}$ containing $X$, using the left endpoint of the unit interval as origin. Write $Y_j=\lfloor2^j(Y-T)\rfloor$ for a vertical origin $T$ fixed below. Let $a=\log_2(1/\lambda)$, $d_0=\log_2 b$ and
\[
 k_n=\lceil d_0 n\rceil,\qquad j_n=\lfloor a n\rfloor .
\]
Take every finitely supported integer vector $h=(h_i)$ and its
ambient character $\chi_h(\theta)=\exp(2\pi i\sum_i h_i\theta_i)$
on $\T^{\mathbb Z}$. Its restriction to $H$ is nontrivial precisely
when $D=\sum_i h_i b^i\ne0$. For each such $h$,
\[
 \chi_h(\theta_{q\ell}(x))
  = \exp\!\left(2\pi i\sum_i h_i\varphi_{q\ell+i}\right)
    \exp(2\pi i D b^{q\ell}x).
\]
The first factor is deterministic and has modulus one. On the unit input
interval, integrating the square of the full character average gives
\[
 \int \left|\frac{1}{N}\sum_{\ell=1}^N \chi_h(\theta_{q\ell}(x))\right|^2
 dx \le \frac{C_{h,b,q}}{N}.
\]
Indeed, the diagonal terms contribute $1/N$. The absolute values of
the off-diagonal integrals are bounded by $1/(\pi|D|\,|b^{qv}-b^{qw}|)$, whose sum over $v>w\ge1$ converges. Chebyshev and Borel--Cantelli along $N=j^2$ followed by interpolation show convergence to zero almost surely. This holds simultaneously for all the countably many characters and $q$. On any interval of scale indices starting and ending at fixed positive multiples of $N$ up to $O(1)$, with length comparable to $N$, the same character averages vanish by subtraction. Every weak limit of the empirical measures on such a window is $H$-translation invariant, since characters nontrivial on $H$ integrate to zero. Its integral of $f_q$ is at least $\inf_\omega\int_H f_q(z+\omega)\,dm_H$ by Haar averaging.
For zero phases take the zero extension to negative indices. Then every
$\theta_n(x)$ belongs to $H$. Every such window limit is supported on
$H$ and invariant under all its translations, so it is Haar probability
on $H$. This holds for Lebesgue-almost every $x$ on the chosen interval.

Evaluate conditional masses at the labels realized at $x$. The horizontal cell at depth $k_n$ has length in $(b^{-n}/2,b^{-n}]$. Matching the vertical depth-$j_{n+q}$ label at a second point $x+b^{-n}u$ in that cell forces $|u|\le1$ and $|W(x+b^{-n}u)-W(x)|\le2\lambda^{n+q}$. Normalizing by the horizontal cell length produces a factor at most two in $du$. For each fixed $q$ and all large $n$ the tail bound gives
\[
 -\log\mathbb P(Y_{j_{n+q}}\mid X_{k_n})\ge f_q(\theta_n(x)).
 \tag{G2}
\]
For the coarser label set $g(t)=-\log\operatorname{dist}(t,\mathbb Z)$. By H\"older continuity an interval inside the horizontal cell of length at least $C b^{-n}\operatorname{dist}(2^{j_n}(Y-T),\mathbb Z)^{1/\alpha}$ stays in the same vertical cell. Hence
\[
 -\log\mathbb P(Y_{j_n}\mid X_{k_n})
  \le C_0+\alpha^{-1}g(2^{j_n}(Y-T)).
 \tag{G3}
\]
The ergodic theorem for doubling applied to $g^2\in L^1(\T)$ and Fubini in $T$ give one $T$ for which, almost surely in $x$, $\sum_{j=0}^K g(2^j(Y-T))^2\le C_1K$ for all large $K$; $C_1$ can be fixed larger than $\int_\T g^2$.

Choose $q$ with $a q>1$ and $d_0 q>1$. At target depth $K$ use the consecutive multiples $n_0,\dots,n_\ell$ of $q$ beginning with the first $k_{n_0}\ge K$ and ending with the last $j_{n_\ell}\le K$. Then
\[
 \ell=K(1/a-1/d_0)/q+O(1).
\]
Let $p_i$ be the conditional law of the next vertical label given $(Y_{j_{n_i}},X_K)$, and $r_i$ its law given $(Y_{j_{n_i}},X_{k_{n_i}})$. The latter fields $\mathcal F_i$ are increasing, and for the realized labels
\[
 \mathbb E(p_i/r_i\mid\mathcal F_i)
   =\sum_{z:r_i(z)>0}p_i(z)\le1 .
\]
The product of preceding ratios is measurable in $\mathcal F_i$. Markov's inequality for the product, followed by Borel--Cantelli with thresholds $\exp(\epsilon K)$ for rational $\epsilon>0$, gives $\sum_i-\log p_i\ge\sum_i-\log r_i-o(K)$ almost surely. The ratio of fine to coarse masses conditional on $X_{k_{n_i}}$ equals the realized $r_i$. By (G2), (G3) and Cauchy--Schwarz,
\[
 \sum_i-\log r_i
 \ge \sum_{i=0}^{\ell-1}f_q(\theta_{n_i}(x))
          -C_0\ell-\alpha^{-1}\sqrt{\ell C_1K}.
\]
Since $j_{n_i}$ are distinct and at most $K$, the boundary loss just
displayed, divided by $K$, has limsup $O(q^{-1}+q^{-1/2})$ at fixed $q$.
The graph mass of the realized depth-$K$ square is at most $2^{-K}\prod_i p_i$ by the conditional chain rule. Apply the character average bound to this window at fixed $q$, divide by $K\log2$, then let $q\to\infty$ and use (G1). The lower dyadic local dimension is at least $1+s(1-\alpha)$ almost everywhere. The strict order of these limits permits all the constants at a fixed block length to depend on it.

For each smaller exponent $t$, a restriction of positive mass to points whose dyadic cell masses are at most $2^{-kt}$ at all sufficiently large depths gives the Hausdorff lower bound, since any ball of comparable radius meets boundedly many cells. This also applies after restricting to any positive-length input interval. The H\"older cover needs at most $C2^{K(2-\alpha)}$ squares of side $2^{-K}$ on a unit interval. The total mass of squares of mass less than $2^{-K(2-\alpha+\epsilon)}$ is summable in $K$. When $s=1$, Borel--Cantelli thus gives equality of dyadic information rates. The mass of points whose square has a neighbor at least $2^{K\epsilon}$ times heavier is at most $8\,2^{-K\epsilon}$ by summing over neighboring square pairs. Comparing balls and neighboring squares now gives exact dimensionality in the Euclidean metric. Restriction of measure to subintervals preserves this almost-everywhere local dimension by differentiation.
\end{proof}



\paragraph{Dependence of the exceptional set on the phases.}
Fix an interior point $x_0$ of a bounded input interval $I$ and choose
$\varphi_j=-b^j x_0\bmod 1$. Every summand attains its maximum at $x_0$.
For $0<|t|\le1/4$, choose $k\ge0$ with $1/(4b)<b^k|t|\le1/4$.
All terms in $W_\varphi(x_0)-W_\varphi(x_0+t)$ are nonnegative. The
$k$th term is at least $\lambda^k(1-\cos(\pi/(2b)))\gg_b |t|^\alpha$.
Together with the H\"older upper bound this gives
$W_\varphi(x_0)-W_\varphi(x_0+t)\asymp |t|^\alpha$.
Consequently the graph measure of the radius-$r$ ball centered at
$(x_0,W_\varphi(x_0))$ is comparable to $r^{1/\alpha}$ as $r\downarrow0$:
the horizontal preimage lies between centered intervals of radii
$c r^{1/\alpha}$ and $C r^{1/\alpha}$. Its local dimension at this
point is $1/\alpha>2-\alpha$. Thus the almost-everywhere conclusion of
Theorem~\ref{thm:main} holds with the phase sequence fixed first.

\section{Information from the finite predecessor choices}\label{app:pred}
For any primitive irreducible polynomial with absolute constant coefficient $p\ge2$ we establish the conditional and anchored estimates of the opening section. Take $\gamma=(b\lambda)^{-1}$ in the application to slopes of the original graph.
\subsection{Dimension of the conditional slope measures}\label{sec:slopes}
Let $b>1$ be real algebraic of degree $d$, with primitive minimal
polynomial $a_0+\cdots+a_dX^d$. Set $p=|a_0|$, $c=1/b$ and take
$0<\gamma<1$. Suppose $p\ge2$. Use Haar probability on the phase
group with this polynomial recurrence. Its central block
$w=(z_0,\ldots,z_{d-1})$ is uniform on $\mathbb T^d$. Given $w$,
one samples a predecessor with $p$ equiprobable branches. The new
state is the block $(z_{-1},z_0,\ldots,z_{d-2})$ and its unconditional
law remains Haar. Iterating this rule gives the conditional past.
Indeed a finite admissible path extends in both directions by solving
the circle equations, and the Haar fibers have equal weights.
The annihilator description follows from Gauss's lemma applied to the
primitive minimal polynomial. In particular the shift and every nonzero
power of it are ergodic: its dual action on nonzero elements of the
additive group $\mathbb Z[b,b^{-1}]$ is multiplication by $b$, with
no finite orbit.

Fix a real lift $\widetilde w=(\widetilde w_0,\ldots,\widetilde w_{d-1})$
of the central block. At a node the digit $i\in\{0,\ldots,p-1\}$
selects the predecessor
\[
 y=\frac{i-\sum_{j=1}^d a_j\widetilde w_{j-1}}{a_0}.
\]
Use $(y,\widetilde w_0,\ldots,\widetilde w_{d-2})$ as the inherited
real lift for the next step (just $(y)$ when $d=1$).
All suffix identities with unchanged digits use that inherited lift.
If $\widetilde w$ is replaced by $\widetilde w+k$ with $k\in\mathbb Z^d$,
the same predecessor modulo one has digit
$i'\equiv i+\sum_{j=1}^d a_j k_{j-1}\pmod p$: the two numerators then
differ by a multiple of $a_0$. Repeating this argument at each child
permutes the digits at that node and preserves every cylinder's
probability and phase values. Thus conditional laws are defined on the
torus. Pathwise derivatives below are computed on the real cover with
the digits fixed. On the real direction
$v=(1,b,\ldots,b^{d-1})$, $z_{-l}(w+tv,I)=z_{-l}(w,I)+c^l t$
modulo one for a fixed lifted code, by the polynomial relation.
Set $g_j=\gamma c^j$ and
\[
 F_j(w,I)=\sum_{l\ge1}g_j^{l-1}\sin(2\pi z_{-l}(w,I)+j\pi/2),
 \qquad \nu^j_w=\operatorname{law}_I F_j(w,I),\quad j\ge0 .
\]
Let $S_j^n$ denote the first $n$ terms. Along $v$, all derivatives
converge uniformly and $\partial_v F_j=2\pi c F_{j+1}$, with the
same identity for truncations. We write $w_n$ for the state reached
after $n$ predecessors. For every infinite continuation $I'$ and
prefix $u$ of length $n$,
\[
      F_j(w,uI')=S_j^n(w,u)+g_j^n F_j(w_n,I').
\]

\begin{theorem}[Conditional slope dimension]\label{thm:algebraic-slope-dimension}
For every $j\ge0$ and almost every $w$, the conditional measure
$\nu^j_w$ is exact dimensional, with dimension
\[
                 s_j=\min\{1,\log p/\log(1/g_j)\}.
\]
\end{theorem}
\begin{proof}
At depth $k$ the support is covered by $p^k$ intervals of length at most
$2g_j^k/(1-g_j)$. The uniform code also puts mass at least $p^{-k}$
in the ball of this radius about each sampled point. Together with the
usual almost-everywhere upper local dimension bound of one on the line,
this gives the required upper bound. We prove the lower bound.

\par\smallskip\noindent\textbf{Initialization.}\enspace
Take $j$ large enough that $g_j/(1-g_j)<\sin(\pi/p)$. For any two
codes differing at the first digit, let $f(w)$ be their difference
of $F_j$ values. In $(f,(2\pi c)^{-1}\partial_v f)$ the leading
vector has norm at least $2\sin(\pi/p)$, by the quadrature of sine
and cosine. The $l$th vector has norm at most $2g_j^{l-1}$, since
its second coordinate carries an additional factor $c^{l-1}$.
Hence the norm of the full pair is bounded away from zero. For all
small $r$, $|f|\le r$ forces $|\partial_v f|\ge\epsilon>0$
uniformly over these pairs.

Here and below we use the bound
\[
 \operatorname{Leb}\{w\in[0,1)^d:|f(w)|\le r,
                       |\partial_v f(w)|>e\}\le C_e r .
\]
For all functions at issue $|\partial_v^2 f|$ is uniformly bounded.
Slice the unit cube into intervals parallel to $v$, using its first
coordinate as the interval parameter. Each has length at most one in
this parameter. Split each interval into boundedly many pieces so small
that the derivative varies by less than $e/2$. On any piece containing
a counted point it has constant sign and modulus at least $e/2$.
The mean value theorem and integration over the bounded set of transverse
parameters give the claim. Real lifts allow this argument without
regularity assumptions in transverse directions.

Two independent codes share exactly $k$ initial digits with probability
$p^{-k}(1-p^{-1})$. After that prefix the difference of the two
functions is $g_j^k$ times a fresh difference at $w_k$. Averaging the
shared code and $w$ restores the uniform law for $w_k$ modulo one.
Lifts may be recoded there while preserving the condition of first
disagreement. The preceding sublevel bound gives
\[
   \int (\nu^j_w\otimes\nu^j_w)\{|x-y|\le r\}\,dw
       \le C\sum_{k\ge0}p^{-k}\min(1,r g_j^{-k}).
\]
This is $O_t(r^t)$ for every $t<s_j$. Thus for almost every $w$ the
$t'$-energy of $\nu^j_w$ is finite for every $0<t'<s_j$, by taking
$t'<t<s_j$ and integrating in the radius. This proves the lower
local dimension bound at these large indices.

\par\smallskip\noindent\textbf{The induction step.}\enspace
Fix $j$. Write $g=g_j$, $h=g_{j+1}$ and suppose that $\nu^{j+1}_w$ has
lower local dimension at least $s>0$ almost everywhere for almost every
$w$. Choose
\[
       0<t<\min\{1,s\log(1/h)/\log(1/g)\}.
\]
For a block length $n$, join two length-$n$ prefixes at $w$ if
\[
   |\partial_v S_j^n(w,u)-\partial_v S_j^n(w,v')|
                     \le (2C+1)h^n ,
\]
where $C$ bounds $h^{-n}|\partial_v(F_j-S_j^n)|$. Count a prefix
as its own neighbor. Distinct full paths whose prefixes are not
neighbors have directional derivative difference greater than $h^n$.
Call a prefix bad if its number of neighbors exceeds $p^n g^{tn}$.
For a sampled infinite path $I$ all extensions of every neighbor
of its prefix have $F_{j+1}$ value within $C' h^n$ of
$F_{j+1}(w,I)$. This follows from the derivative identity for
truncations and the uniform tail bound. Thus, pointwise, the degree
divided by $p^n$ is at most
$\nu^{j+1}_w([F_{j+1}(w,I)-C'h^n,F_{j+1}(w,I)+C'h^n])$.
The local dimension hypothesis and the strict inequality defining $t$
show that the Haar probability $e_n$ of a bad prefix tends to zero.

\par\smallskip\noindent\textbf{Correlation estimate for a prescribed sequence of blocks.}\enspace
For a binary word $P$ of length $k$, let $\mu_{w,P}$ be the
unnormalized pushforward under $F_j(w,\cdot)$ of paths with those
successive good or bad statuses in blocks of length $n$. Marking a
block only uses the state and the digits in that block. For $r>0$ set
$U_r=\mathrm{1}_{[0,r]}/r$ and
\[
       E_P(r)=\int\|U_r*\mu_{w,P}\|_2^2\,dw .
\]
The correlation kernel $K_r=U_r*U_r(-\cdot)$ equals
$r^{-1}(1-|x|/r)_+$. Write $P'$ for $P$ without its initial status.
We claim
\[
 E_P(r)\le C_n+A E_{P'}(r/g^n),\qquad
 A=\begin{cases}g^{-n(1-t)},&\text{initial status good},\\
                  g^{-n},&\text{initial status bad}.
    \end{cases}                                                     \tag{*}
\]
Expand the squared norm according to the two first prefixes. For
pairs which are not neighbors at $w$, discard all path restrictions
and retain the necessary condition that the full directional
derivative difference is greater than $h^n$. At fixed two full
codes, the directional sublevel estimate bounds the integral of
$K_r(F_j(w,I)-F_j(w,I'))$ on that set by $C_n$, uniformly in the
codes and $r$. Integrating the independent codes proves the first
term in $(*)$.

For neighboring pairs, at fixed $w$, bound each scalar product
of the smoothed child measures by half the sum of their squared
norms. Both children have the prescribed status. Each diagonal
norm therefore occurs with multiplicity at most $p^n g^{tn}$ in
the good case and at most $p^n$ in the bad case. Each pair has
weight $p^{-2n}$. In a diagonal norm, translating the prefix and
scaling by $g^n$ produce $g^{-n}\|U_{r/g^n}*\mu_{w_n,P'}\|_2^2$.
Drop the initial status restriction and sum all prefixes. Their
uniform weights $p^{-n}$ together with $dw$ give Haar $dw_n$.
The status of the remaining blocks respects relabeling when the lift
of $w_n$ is changed. This proves $(*)$.

\par\smallskip\noindent\textbf{From typical block sequences to local dimension.}\enspace
Fix $n$ sufficiently large that $e_n<1/2$ and choose $e_n<\eta<1/2$.
By ergodicity of the shift by $n$, almost every path has at most
$\eta k$ bad blocks among its first $k$ blocks for every sufficiently
large $k$. For $k_0\ge1$ restrict the original path measure to paths
meeting this condition simultaneously for all $k\ge k_0$; write
$\mu_{w,k_0}^*$ for its vertical pushforward at $w$. At $r_k=g^{nk}$
and for a status word with at most $\eta k$ bad blocks, iteration
of $(*)$ yields
\[
  E_P(2r_k)\le C_n(k+1)g^{-nk(1-t+t\eta)} .
\]
Indeed all coefficients $A$ in $(*)$ are at least one and the
terminal energy at radius $2$ is bounded. There are at most
$\exp(k H(\eta))$ such words, where
$H(\eta)=-\eta\log\eta-(1-\eta)\log(1-\eta)$. For $k\ge k_0$ the
restricted measure is dominated by the sum over these words. The
triangle inequality in $L^2(dw\,dy)$ bounds its squared norm by the
preceding bound times $\exp(2kH(\eta))$. The kernel $K_{2r_k}$ is
at least $1/(4r_k)$ on $[-r_k,r_k]$. Consequently
\[
 \int (\mu^*_{w,k_0}\otimes\mu^*_{w,k_0})\{|x-y|\le r_k\}\,dw
 \le 4 C_n(k+1)r_k^{t(1-\eta)-2H(\eta)/(n\log(1/g))}.
\]
Since the $r_k$ form a geometric progression, for almost every $w$
this measure has finite $s'$-energy for each positive $s'$ strictly
below the exponent displayed. Its lower local dimension is at least
$s'$ almost everywhere. For almost every $w$, the measures
$\mu^*_{w,k_0}$ increase to $\nu^j_w$ by the ergodic conclusion and
disintegration. Differentiation of these submeasures relative to
$\nu^j_w$ transfers their almost-everywhere lower bounds to
$\nu^j_w$: almost every point has positive Radon--Nikodym derivative
for some $k_0$. As $n\to\infty$ one can choose $\eta\downarrow0$.
This proves the lower bound $t$. Taking a countable sequence of such
$t$ proves
\[
  \underline{\dim}_{\mathrm{loc}}\nu^j_w
       \ge \min\{1,s\log(1/g_{j+1})/\log(1/g_j)\}
       \quad\text{almost everywhere}.
\]
Initialize at a sufficiently large derivative index and iterate
downwards. At each step $g_{j+1}<g_j$ and the asserted $s_j$ follows.
\end{proof}

For the classical two-sided increment one has
$\Theta_0=2\pi\gamma F_0$ and $\gamma=1/(b\lambda)$.
Thus the conditional slope has dimension one whenever $p\gamma\ge1$.
In particular, every rational $b=P/Q>1$ in lowest terms satisfies
this for every $1/b<\lambda<1$, since $p\gamma=Q/\lambda>1$.

\par\smallskip\noindent\textbf{Uniformity under arbitrary deterministic phases.}\enspace
For a phase string $\phi=(\phi_l)_{l\ge1}$, replace every
$2\pi z_{-l}$ by $2\pi(z_{-l}+\phi_l)$ throughout and denote the
resulting laws by $\nu^{j,\phi}_w$. For every $j$, every $s'<s_j$
and each fixed $A>0$, one has
\begin{equation}\label{eq:uniform-slope-information}
 \lim_{r\downarrow0}\sup_\phi\int
  \nu^{j,\phi}_w\{y:\nu^{j,\phi}_w([y-Ar,y+Ar])>r^{s'}\}\,dw=0 .
\end{equation}
To prove this, use the same induction. At high derivative index the
first-digit quadrature and sublevel bounds are uniform in $\phi$.
The common-prefix conditioning changes the remaining phase string
to its deterministic suffix. The resulting integrated pair bound
and Markov's inequality give \eqref{eq:uniform-slope-information}.
At a lower derivative index the induction hypothesis, applied to
$C'h^n$ as above with an exponent strictly between
$t\log(1/g)/\log(1/h)$ and $s$, makes the bad-block probability
at most $e_n\to0$, uniformly over all starting phase strings.
Recursion $(*)$ is unchanged with the suffix of $\phi$ in the
remaining energy. Its sliced sublevel constant does not depend on
$\phi$. Every $n$-block therefore has bad probability at most $e_n$,
irrespective of the block position, by stationarity of the underlying
Haar states.

For given $k$ write $\widetilde\mu_w$ for the submeasure of paths
having at most $\eta k$ bad blocks among their first $k$. The
discarded mass, averaged over $w$, is at most $e_n/\eta$ by Markov's
inequality. Set $r_k=g^{nk}$ and
$t_*=t(1-\eta)-2H(\eta)/(n\log(1/g))$. The prescribed-word
estimate gives
\[
 \int(\widetilde\mu_w\otimes\widetilde\mu_w)
       \{|x-y|\le A' r_k\}\,dw\le C_{n,A'}(k+1)r_k^{t_*}
\]
for each fixed $A'>0$, by starting the recursion at $2A'r_k$.
For each $w$, apply Lemma~\ref{lem:full-ball} with
$\nu=\nu^{j,\phi}_w$, retained measure $\widetilde\mu_w$, $R=A r_k$
and $\tau=r_k^{s'}$. Averaging over $w$ bounds the probability in
\eqref{eq:uniform-slope-information} by
\[
 e_n/\eta+3/M+M C_{n,A}(k+1)r_k^{t_*-s'} .
\]
Thus if $s'<t_*$ the limsup of
the probability in \eqref{eq:uniform-slope-information} on this
scale sequence is at most $e_n/\eta+3/M$. The same conclusion holds
between consecutive scales by increasing the fixed radius constant
by $g^{-n}$. Here $n$ stays fixed as $k\to\infty$. Now choose $t$
strictly between the desired $s'$ and the upper endpoint in the
induction step. Send $n\to\infty$ with $\eta\to0$ and
$e_n/\eta\to0$, and then send $M\to\infty$. This proves the
uniform information bound. The changing submeasure in this argument
is used only at its own scale.

In Theorem~\ref{thm:conditional-graph} the central block can be
$w\in\mathbb T^d$ throughout. The proof tests the conditional mass
only at finitely many radii $C'g_0^{d_k}$ with $d_k\to\infty$.
The vanishing exceptional probability there follows from
\eqref{eq:uniform-slope-information}, uniformly for the phase strings
at the different starting depths. Prefix sums carry the appropriate
suffix phases, and the exact increment identity holds along $v$.
Thus the lower bound for the integral of the entropy cutoff holds
uniformly on all translates of the Haar phase group. For each fixed
deterministic two-sided phase sequence, for Lebesgue-almost every
physical starting point $x$, every weak limit of the empirical measures
of $\theta_n(x)$ on the scale windows of Appendix~\ref{app:graph}
is invariant under translation by $H$. Indeed, each character nontrivial
on $H$ has zero time-average by the mean-square lacunary estimate
with coefficients of modulus one. Such a limit is a mixture of
Haar measures on translates of $H$, by averaging continuous
functions over $H$. The same lower bound for every coset therefore
applies to each weak limit. This gives the graph lower bound
$1+s_0(1-\alpha)$ for each deterministic choice of phases and
the full graph and graph-measure conclusion when $p\gamma\ge1$.

\subsection{From conditional slopes to secant probabilities}\label{sec:secant}
We use $c=1/b$, $\gamma=c/\lambda$, and the Haar environment of the
preceding section. Write $d_b=[\mathbb Q(b):\mathbb Q]$ for the degree.
Given the central block $w\in\T^{d_b}$, any past prefix
of length $n$ and any future prefix of length $m$ extend jointly:
the recurrence uses $d_b+1$ consecutive coordinates, so the central
block supplies all boundary data on both sides. There are
$p^n|a_{d_b}|^m$ such pairs. Projection of conditional Haar measure to
this finite fiber gives equal mass $p^{-n}|a_{d_b}|^{-m}$ to every pair.
This proves conditional independence for the two infinite branches
by their cylinders. The block $w_n$ reached after $n$ past steps has
Haar law by shift invariance; $z_{-n}$ denotes its first coordinate.
The formulas below display zero phases. For prescribed phases add the
corresponding shift to each sine and use its suffix at every new block.
For $|u|\le1$, write the contribution of the negative indices to
the two-sided increment as $-u\Theta_u(w,I)$, where
\[
 \Theta_u=2\pi\sum_{j\ge1}\gamma^j\operatorname{sinc}(\pi c^ju)
          \sin\bigl(2\pi(z_{-j}+c^ju/2)\bigr),
 \qquad \operatorname{sinc}(v)=\sin(v)/v,\quad \operatorname{sinc}(0)=1 .
\]
Write $T_u^n$ for its first $n$ terms. There is a uniform $C_0\ge1$ with
\[
 |\Theta_u|\le C_0,\qquad
 |\Theta_u-T_u^d|\le C_0\gamma^d,\qquad
 |T_u^d-T_0^d|\le C_0|u|\quad(d\ge0).
\]
The third inequality follows by summing $O(\gamma^jc^j|u|)$.
At any prefix one has the exact identity
\[
 T_u^{n+d}(w,I)=T_u^n(w,I)+
     \gamma^n T_{c^nu}^d(w_n,\sigma^n I).
\]
Denote by $\nu_w$ the conditional law of $\Theta_0(w,I)$.

\begin{theorem}[Conditional slopes and graph dimension]\label{thm:conditional-graph}
Suppose $0<s\le1$ and for almost every $w$,
\[
 \liminf_{a\downarrow0}
 \frac{\log\nu_w([y-a,y+a])}{\log a}\ \ge s
 \qquad\text{for $\nu_w$-almost every $y$.}
\]
For the classical zero-phase cosine series at base $b$, the graph
on every interval of positive length has Hausdorff dimension at least
$1+s(1-\alpha)$, where $\alpha=-\log\lambda/\log b$. At $s=1$ the
dimension equals $2-\alpha$ and the graph lift of Lebesgue measure
on each bounded interval is exact dimensional.
\end{theorem}
\begin{proof}
Fix $0<\sigma<s$, and choose $0<\kappa=1/K<1/2$ with $K$ an integer.
All constants in the proof may depend on these choices.
Put $a_0=\log\gamma/\log(c\gamma)\in(0,1)$ and take $L$ fixed
with $a_0^L\le\kappa$. For a small target radius $r>0$, divide
$r^{1-\kappa}<|u|\le1$ into the $K-1$ bands
\[
    r^{(h+1)\kappa}<|u|\le r^{h\kappa},\qquad 0\le h\le K-2.
\]
In band $h$ write $\Delta=r^{1-(h+1)\kappa}$. Thus $r/|u|\le\Delta$.
Choose $n_L=\lfloor\log\Delta/\log\gamma\rfloor$ and define the
preceding depths by $n_{j-1}=\lceil a_0 n_j\rceil$. As $r\downarrow0$
all $n_j$ and every difference $d_j=n_j-n_{j-1}$ tend to infinity
in each band. At $n=n_{j-1}$, $d=d_j$, we have
\[
     c^n\le\gamma^d,\qquad
     \Delta\le\gamma^{n_L}\le\gamma^{-1}\Delta,\qquad
     n_0\le\kappa n_L+(1-a_0)^{-1}.
\]

Choose $A=1+C_0$ and a fixed large $D$, for instance $D=5+8C_0$.
Call a length-$(n+d)$ prefix good at step $(n,d)$ when
\[
 \nu_{w_n}\!\left([T_0^d(w_n,\sigma^nI)-D\gamma^d,
                         T_0^d(w_n,\sigma^nI)+D\gamma^d]\right)
                 \le \gamma^{\sigma d}.
\]
This condition is determined by the prefix, including $w$. For a random
prefix the center here is within $C_0\gamma^d$ of
$\Theta_0(w_n,\sigma^nI)$. The joint law of
$(w_n,\Theta_0(w_n,\sigma^nI))$ is $dw\,d\nu_w(y)$.
In particular, for every starting depth $n$ the probability of a bad
prefix is at most
\[
 \int_{\mathbb T^{d_b}}\!\int
 {\bf 1}_{\{\nu_w([y-(D+C_0)\gamma^d,y+(D+C_0)\gamma^d])
                         >\gamma^{\sigma d}\}}\,d\nu_w(y)\,dw .
\]
The integral is independent of $n$. Its integrand tends to zero at
almost every fixed pair $(w,y)$ by the dimension hypothesis and
$\sigma<s$. Dominated convergence therefore gives a probability tending
to zero even when $n$ varies with $d$. If $E_r$ denotes the union of the bad
events over all bands and their steps, $m_H(E_r)=o(1)$.
There are only $(K-1)L$ such events, and $E_r$ depends solely
on $w$ and its past.

Fix $w$, a future branch code, $u$ in the chosen band, and a target
$y$ depending on these fixed variables. The event
$|\Theta_u-y|\le r/|u|$ forces all prefix windows
\[
       |T_u^{n_j}-y|\le A\gamma^{n_j},\qquad 0\le j\le L .
\]
Now fix the past through depth $n=n_{j-1}$. Write
$z=(y-T_u^n)/\gamma^n$. For every child prefix satisfying the next
window, $|T_{c^nu}^d-z|\le A\gamma^d$, so
$|T_0^d-z|\le(A+C_0)\gamma^d$. For every infinite extension of
each such child,
$|\Theta_0(w_n,\sigma^n I)-z|\le(A+2C_0)\gamma^d$.
If any such child is good, its center $T_0^d$ therefore has a ball
of radius $D\gamma^d$ containing the last interval; here it suffices
that $D\ge 2A+3C_0$. The conditional probability of the union of all
good children satisfying that window is at most $\gamma^{\sigma d}$,
by the definition of that good child and the conditional law $\nu_{w_n}$.
If there is no such good child, the conditional probability is zero.
Both the good tests and the windows are measurable at their respective
prefix depths. Iteration thus bounds the conditional probability of
satisfying all the windows and good tests by
\[
        \gamma^{\sigma(n_L-n_0)}
                  \le C\Delta^{\sigma(1-\kappa)}.
\]
In particular this bound applies after deleting $E_r$.

Conditional on the independent future code, the nonnegative-index
contribution to $V(z,u)$ is a function of $w,u$ and that code.
We can use its quotient by $u$ as the target $y$, then average over
$w$ and the code. Integration on band $h$ multiplies the last
estimate by at most $2r^{h\kappa}$. Its exponent of $r$ is
\[
    h\kappa+(1-(h+1)\kappa)\sigma(1-\kappa)
           \ge\sigma(1-\kappa)^2.
\]
The omitted interval $|u|\le r^{1-\kappa}$ has length $2r^{1-\kappa}$.
Together this gives
\[
    \int_{H\setminus E_r}\int_{-1}^1
           \mathrm{1}_{\{|V(z,u)|\le r\}}\,du\,dm_H(z)
       \le C r^{\sigma(1-\kappa)^2},\qquad m_H(E_r)=o(1).
\]
In the graph entropy estimate take the ceiling for horizontal dyadic
depths, so that normalized shifts obey $|u|\le1$, and take $r=4\lambda^k$.
The continuous cutoff $B_k$ of Appendix~\ref{app:graph} is dominated by twice the
integral in $u$ above. Jensen's inequality on $E_r^c$ gives
\[
 \int_H-\log\min(1,B_k)\,dm_H
     \ge(1-o(1))\bigl[\sigma(1-\kappa)^2
                    k\log(1/\lambda)-O(1)\bigr].
\]
For zero phases and Lebesgue-almost every $x$ on the input interval,
the empirical laws of $\theta_n(x)$ converge to $m_H$ on the windows
used in Appendix~\ref{app:graph}. The point $x=0$, for example, has
$\theta_n(0)=0$ at every scale when the phases and their extension are
zero. The graph estimate uses the almost-everywhere statement because
the natural graph measure is the image of Lebesgue measure.
The moving-partition argument gives graph dimension at least
$1+\sigma(1-\kappa)^2(1-\alpha)$. Let $\sigma\uparrow s$ and
$\kappa\downarrow0$. For $s=1$ the upper covering estimate also
gives exact dimensionality as in Proposition~\ref{prop:graph}.
\end{proof}

If the conditional laws have densities $h_w$ with
$\int\|h_w\|_2^2\,dw<\infty$, the hypotheses hold with $s=1$.
Moreover the probability of a bad prefix is
$O(\gamma^{(1-\sigma)d})$, since the expected mass of a ball of
radius $(D+C_0)\gamma^d$ at a sampled point is at most
$2(D+C_0)\gamma^d\int\|h_w\|_2^2\,dw$ by Cauchy--Schwarz.
Each fixed-band depth difference $d$ is bounded below by a positive
multiple of $\log(1/r)$ for small $r$, so $m_H(E_r)=O(r^\delta)$
for some $\delta>0$. Absolute continuity of $\nu_w$ alone already
implies the qualitative full-dimension criterion by differentiation.


For an arbitrary deterministic phase coset the finite-band argument uses the same conditional past tree. In each stage use the phase-shifted slope law at the new state. The uniform sampled mass estimate above makes the finitely many discarded events have total mass $o(1)$ uniformly in the coset. Given the central coordinates, the future remains independent of the past. On the retained set the averaged anchored mass is therefore $O(r^{\sigma(1-\kappa)^2})$. For $s'<\sigma(1-\kappa)^2$, Markov gives a probability at most $o(1)+O(r^{\sigma(1-\kappa)^2-s'})$ that the full anchored mass exceeds $r^{s'}$. For arbitrary fixed $T,A>0$, apply \eqref{eq:fixed-dilation} with
$B=b$, $v=1$ and $c_j=1$. Taking $\sigma$ up to the conditional
exponent and $\kappa$ down to zero gives $\mathcal S_{s_0}$ with $s_0=\min(1,\log p/\log(1/\gamma))$ as in Theorem~\ref{thm:algebraic-slope-dimension}.


\section*{Manuscript preparation}
The author used GPT (OpenAI) and Claude (Anthropic) to test conjectures,
derive and check formulas, draft arguments, and check steps in the proposed
proof. Both assistants also contributed to the ongoing Lean~4 formalization
of parts of the argument. Verification of the full theorem in Lean remains
incomplete.

\begin{thebibliography}{10}
\bibitem{mw} R. D. Mauldin and S. C. Williams, \href{https://sites.itservices.cas.unt.edu/~mauldin/papers/no60.pdf}{On the Hausdorff dimension of some graphs}, Trans. Amer. Math. Soc. 298 (1986), 793--803.
\bibitem{hunt} B. R. Hunt, \href{https://doi.org/10.1090/S0002-9939-98-04387-1}{The Hausdorff dimension of graphs of Weierstrass functions}, Proc. Amer. Math. Soc. 126 (1998), 791--800.
\href{https://user.eng.umd.edu/~bhunt/papers/hdimw.pdf}{Author's PDF}.
\bibitem{bbr} K. Bara\'nski, B. B\'ar\'any and J. Romanowska,
\href{https://doi.org/10.1016/j.aim.2014.07.033}{On the dimension of
the graph of the classical Weierstrass function}, Adv. Math. 265
(2014), 32--59.
\href{https://arxiv.org/abs/1309.3759}{arXiv:1309.3759}.
\bibitem{shen} W. Shen, \href{https://arxiv.org/abs/1505.03986}{Hausdorff dimension of the graphs of the classical Weierstrass functions}, Math. Z. 289 (2018), 223--266.
\bibitem{keller} G. Keller, \href{https://doi.org/10.1214/15-AIHP711}{A simpler proof for the dimension of the graph of the classical Weierstrass function}, Ann. Inst. Henri Poincar\'e Probab. Stat. 53 (2017), 169--181.
\href{https://arxiv.org/abs/1406.3571}{arXiv:1406.3571}.
\bibitem{ren} H. Ren, \href{https://doi.org/10.1016/j.aim.2024.109587}{A dichotomy for the dimension of SRB measure}, Adv. Math. 442 (2024), article 109587.
\href{https://arxiv.org/abs/2208.04576}{arXiv:2208.04576}.
\bibitem{rs} H. Ren and W. Shen,
\href{https://doi.org/10.1007/s00222-021-01060-2}{A dichotomy for the
Weierstrass-type functions}, Invent. Math. 226 (2021), 1057--1100.
\href{https://arxiv.org/abs/2007.04312}{arXiv:2007.04312}.
\bibitem{lptx} J. Li, W. Pan, Y. Tong and D. Xu, \href{https://arxiv.org/abs/2608.30744}{Dimensions of surface repellers and attractors of non-linear planar IFSs}, preprint (2026).
\bibitem{ledrappier} F. Ledrappier, \href{https://doi.org/10.1090/conm/135/1185095}{On the dimension of some graphs}, in \emph{Symbolic Dynamics and its Applications} (P. Walters, ed.), Contemp. Math. 135 (1992), 285--293.
\bibitem{rz} H. Ruan and N. Zhang, \href{https://doi.org/10.4208/ata.OA-SU8}{Hausdorff dimension of a class of Weierstrass functions}, Anal. Theory Appl. 36 (2020), 482--496.
\href{https://www.global-sci.com/ata/article/view/8206}{Journal page}.
\end{thebibliography}
\end{document}
